\setgebiet{7}{Geometrie}{mittel}
\chapter{Lösungen: Geometrie}

%%%%%%%%%% 6.1
\stupserzwei{tief}{(a) Die Dreiecke $SAP$ und $SBP$ haben $SA = SB$, $AP = BP$ und die gemeinsame Seite $SP$: SSS. Also sind die Winkel bei $S$ gleich.}{22pt}
\begin{loesungkarte}{tief}{6.1}{Winkelhalbierende und Lot}
\derweg{Konstruieren, dann zwei Dreiecke finden, die nach einem Kongruenzsatz gleich sind --- die gesuchte Gleichheit ist eine Folge davon.}
\schritt{(a) Konstruktion}
Kreis um $S$ mit beliebigem Radius schneidet die Schenkel in $A$ und $B$. Kreise um $A$ und $B$ mit gleichem Radius $> \frac12AB$ schneiden sich in $P$ (im Winkelfeld). Die Gerade $SP$ ist die Winkelhalbierende.
\schritt{(a) Beweis}
In den Dreiecken $SAP$ und $SBP$ gilt $SA = SB$ (gleicher Kreis um $S$), $AP = BP$ (gleiche Radien), $SP$ gemeinsam. Nach SSS sind sie kongruent, also $\angle ASP = \angle BSP$. $\blacksquare$
\schritt{(b) Konstruktion}
Kreis um $P$, der $g$ in $A$ und $B$ schneidet. Kreise um $A$ und $B$ mit gleichem Radius schneiden sich in $Q$ (auf der anderen Seite von $g$). $PQ$ ist das Lot.
\schritt{(b) Beweis}
$PQ$ ist die Mittelsenkrechte von $AB$ (Aufgabe 6.3): $P$ und $Q$ haben beide gleichen Abstand zu $A$ und $B$. Also steht $PQ$ senkrecht auf $AB = g$. $\blacksquare$
\schritt{(c)}
Halbieren braucht nur Kreise gleicher Radien --- eine quadratische Operation (Schnitt zweier Kreise). Dritteln des Winkels $60^{\circ}$ hieße $\cos20^{\circ}$ konstruieren, und das ist Nullstelle von $8x^{3} - 6x - 1$, eines irreduziblen kubischen Polynoms. Mit Zirkel und Lineal erreicht man nur Zahlen, deren Grad eine Zweierpotenz ist (Wantzel 1837, Band 11).
\probe{(a) Probe am rechten Winkel: $A = (1, 0)$, $B = (0, 1)$, $P = (1, 1)$ (Schnitt der Kreise mit Radius $1$); $SP$ ist die Diagonale $y = x$, Winkel $45^{\circ}$. \checkmark}
\begin{mitnehmen}\kastentext
Jede Konstruktion ist ein Satz: \enquote{Wenn man so zeichnet, dann gilt \dots} Der Beweis benennt die kongruenten Dreiecke und den Kongruenzsatz. Und fast immer ist es SSS, weil der Zirkel gleiche Längen liefert.
\end{mitnehmen}
\end{loesungkarte}

%%%%%%%%%% 6.2
\stupserzwei{akzent}{$\sin\alpha = \frac{a\sin\beta}b = \frac{5\cdot0{,}5}3 = 0{,}833$: $\alpha = 56{,}4^{\circ}$ oder $123{,}6^{\circ}$. Beide passen, weil $\alpha + \beta < 180^{\circ}$ in beiden Fällen.}{22pt}
\begin{loesungkarte}{akzent}{6.2}{SSW ist kein Kongruenzsatz}
\derweg{Konstruieren, mit dem Sinussatz beide Lösungen finden, die Eindeutigkeitsbedingung ablesen.}
\schritt{(a)}
Zeichne $a = BC = 5$. Trage in $B$ den Winkel $\beta = 30^{\circ}$ an. Kreis um $C$ mit Radius $b = 3$: Er schneidet den freien Schenkel zweimal, in $A_1$ und $A_2$ --- weil $b = 3$ größer als der Abstand von $C$ zum Schenkel ($5\sin30^{\circ} = 2{,}5$) und kleiner als $a = 5$ ist. Zwei Dreiecke mit denselben Daten.
\schritt{(b)}
Sinussatz: $\frac{\sin\alpha}a = \frac{\sin\beta}b$, $\sin\alpha = \frac{5\cdot0{,}5}3 = 0{,}8333$. $\alpha_1 = 56{,}4^{\circ}$, $\alpha_2 = 180^{\circ} - 56{,}4^{\circ} = 123{,}6^{\circ}$. Dann $\gamma_1 = 93{,}6^{\circ}$, $c_1 = \frac{b\sin\gamma_1}{\sin\beta} = 5{,}99$; $\gamma_2 = 26{,}4^{\circ}$, $c_2 = 2{,}67$.
\schritt{(c)}
SsW: Der gegebene Winkel liegt der \emph{größeren} Seite gegenüber. Dann ist $\alpha$ gegenüber der kleineren Seite spitz (der größere Winkel liegt der größeren Seite gegenüber), und die zweite Lösung $180^{\circ} - \alpha$ fällt weg. Hier läge $\beta$ gegenüber $b = 3 < a = 5$: nicht erfüllt. Mit $b = 6 > a$ wäre der Kreis groß genug, nur einmal zu schneiden.
\schritt{(d)}
\enquote{Also sind sie kongruent} --- der Schluss von SSW auf Kongruenz ist falsch; (a) ist das Gegenbeispiel. Richtig wäre er, wenn zusätzlich $b\ge a$ bekannt wäre.
\begin{fehler}\kastentext
Vier Kongruenzsätze: SSS, SWS, WSW (und WWS), SsW. Nicht SSW, nicht WWW (das ist Ähnlichkeit). Der Sinussatz liefert aus $\sin\alpha$ immer zwei Winkel; welcher gilt, entscheidet die Geometrie --- und manchmal beide.
\end{fehler}
\begin{mitnehmen}\kastentext
Bei jeder Dreiecksberechnung mit Sinussatz: Prüfen, ob $180^{\circ} - \alpha$ auch passt. Der Kosinussatz hat dieses Problem nicht --- $\cos$ ist auf $[0^{\circ}, 180^{\circ}]$ injektiv, deshalb liefert SWS über den Kosinussatz immer eindeutig.
\end{mitnehmen}
\end{loesungkarte}

%%%%%%%%%% 6.3
\stupserzwei{gruen}{$APB$ und $AQB$ sind kongruent (SSS), also $\angle PAB = \angle QAB$. Damit sind $APM$ und $AQM$ kongruent (SWS), also $\angle AMP = \angle AMQ$ --- Nebenwinkel, die gleich sind, sind rechte.}{22pt}
\begin{loesungkarte}{gruen}{6.3}{Die Mittelsenkrechte}
\derweg{Zweimal Kongruenz: erst zeigen, dass $PQ$ die Strecke $AB$ halbiert, dann dass die Winkel bei $M$ rechte sind.}
\schritt{Konstruktion}
Kreise um $A$ und $B$ mit gleichem Radius $r > \frac12AB$ schneiden sich in $P$ und $Q$. Gerade $PQ$, Schnitt mit $AB$ in $M$.
\schritt{$PQ$ halbiert $AB$}
Dreiecke $APQ$ und $BPQ$: $AP = BP = r$, $AQ = BQ = r$, $PQ$ gemeinsam: SSS, kongruent. Also $\angle APQ = \angle BPQ$. Nun $APM$ und $BPM$: $AP = BP$, $PM$ gemeinsam, eingeschlossene Winkel gleich: SWS, kongruent. Also $AM = BM$. $M$ ist der Mittelpunkt.
\schritt{$PQ$ steht senkrecht}
Aus derselben Kongruenz: $\angle AMP = \angle BMP$. Die beiden sind Nebenwinkel (sie ergänzen sich zu $180^{\circ}$) und gleich, also je $90^{\circ}$. $\blacksquare$
\schritt{Umkehrung}
Jeder Punkt der Mittelsenkrechten hat gleichen Abstand zu $A$ und $B$, und nur diese Punkte: Die Mittelsenkrechte ist der Ort aller Punkte mit $XA = XB$ --- der Grund, warum sich drei Mittelsenkrechten in einem Punkt treffen (Aufgabe 6.5).
\probe{Koordinaten: $A = (-1, 0)$, $B = (1, 0)$, $r = 2$: Kreise $(x\pm1)^{2} + y^{2} = 4$ schneiden sich bei $x = 0$, $y = \pm\sqrt3$. $PQ$ ist die $y$-Achse. \checkmark}
\begin{mitnehmen}\kastentext
Der Beweis braucht $r > \frac12AB$ (sonst schneiden sich die Kreise nicht) und die Tatsache, dass gleiche Nebenwinkel rechte sind. Das sind genau die Stellen, an denen Euklid in den \emph{Elementen} (I.10 und I.11) ebenfalls argumentiert.
\end{mitnehmen}
\end{loesungkarte}

%%%%%%%%%% 6.4
\stupserzwei{tief}{Die Höhe $12$ auf $b = 14$ teilt sie in $5$ und $9$: Die beiden Teildreiecke sind $5$-$12$-$13$ und $9$-$12$-$15$ --- beide pythagoreisch.}{22pt}
\begin{loesungkarte}{tief}{6.4}{Das 13-14-15-Dreieck}
\derweg{Heron, dann Höhe aus der Fläche, dann $r = A/s$, $R = abc/4A$, dann Kosinussatz als Probe.}
\schritt{(a)}
$s = \frac{13 + 14 + 15}2 = 21$. $A = \sqrt{21\cdot8\cdot7\cdot6} = \sqrt{7056} = 84$.
\schritt{(b)}
$h_b = \frac{2A}b = \frac{168}{14} = 12$. Der Fußpunkt teilt $b$: $\sqrt{13^{2} - 12^{2}} = 5$ und $\sqrt{15^{2} - 12^{2}} = 9$, $5 + 9 = 14$. \checkmark\ Das Dreieck ist aus einem $5$-$12$-$13$- und einem $9$-$12$-$15$-Dreieck zusammengesetzt --- deshalb die schönen Zahlen.
\schritt{(c)}
$r = \frac As = \frac{84}{21} = 4$. $R = \frac{abc}{4A} = \frac{2730}{336} = \frac{65}8 = 8{,}125$.
\schritt{(d)}
$\cos\alpha = \frac{b^{2} + c^{2} - a^{2}}{2bc} = \frac{196 + 225 - 169}{420} = \frac{252}{420} = 0{,}6$, $\alpha = 53{,}13^{\circ}$, $\sin\alpha = 0{,}8$. $A = \frac12\cdot14\cdot15\cdot0{,}8 = 84$. \checkmark
\probe{$\cos\beta = \frac{169 + 225 - 196}{390} = \frac{198}{390} = 0{,}5077$, $\beta = 59{,}5^{\circ}$; $\cos\gamma = \frac{169 + 196 - 225}{364} = \frac{140}{364} = 0{,}3846$, $\gamma = 67{,}4^{\circ}$. Summe $53{,}1 + 59{,}5 + 67{,}4 = 180{,}0$. \checkmark\ Euler: $R\ge2r$: $8{,}125\ge8$. \checkmark}
\begin{mitnehmen}\kastentext
Die Formeln, die man braucht: Heron $A = \sqrt{s(s - a)(s - b)(s - c)}$, $A = rs$, $A = \frac{abc}{4R}$, $A = \frac12bc\sin\alpha$. Vier Ausdrücke für dieselbe Fläche --- und jeder löst eine andere Aufgabe. Das $13$-$14$-$15$-Dreieck ist das kleinste mit ganzzahligen Seiten, Fläche, Höhe und Inkreisradius ohne rechten Winkel.
\end{mitnehmen}
\end{loesungkarte}

%%%%%%%%%% 6.5
\stupserzwei{gruen}{$U = (3, 1)$, $H = (2, 2)$, $S = (\frac83, \frac43)$. Prüfen: $S - U = (-\frac13, \frac13)$, $H - U = (-1, 1)$ --- parallel, Verhältnis $1 : 3$.}{22pt}
\begin{loesungkarte}{gruen}{6.5}{Die Eulergerade}
\derweg{Drei Punkte ausrechnen, dann prüfen, ob die Verbindungsvektoren parallel sind.}
\schritt{Schwerpunkt}
$S = \frac{A + B + C}3 = \bigl(\frac83, \frac43\bigr)$.
\schritt{Umkreismittelpunkt}
Mittelsenkrechte von $AB$: $x = 3$. Mittelsenkrechte von $AC$: Mittelpunkt $(1, 2)$, Richtung von $AC$ ist $(2, 4)$, also $2(x - 1) + 4(y - 2) = 0$, $x + 2y = 5$. Schnitt: $x = 3$, $y = 1$. $U = (3, 1)$. Probe: $UA = \sqrt{9 + 1} = \sqrt{10}$, $UB = \sqrt{9 + 1}$, $UC = \sqrt{1 + 9}$. \checkmark
\schritt{Höhenschnittpunkt}
Höhe von $C$ auf $AB$: $x = 2$. Höhe von $A$ auf $BC$: $BC = (-4, 4)$, also Richtung $(1, 1)$ durch $A$: $y = x$. Schnitt: $H = (2, 2)$. Probe: Höhe von $B$ auf $AC$: Richtung senkrecht zu $(2, 4)$, also $(2, -1)$: $(6, 0) + t(2, -1)$ trifft $(2, 2)$ bei $t = -2$. \checkmark
\schritt{Eulergerade}
$H - U = (-1, 1)$, $S - U = \bigl(-\frac13, \frac13\bigr) = \frac13(H - U)$. $S$ liegt auf $UH$ und teilt die Strecke im Verhältnis $US : SH = 1 : 2$. Genau Eulers Satz.
\probe{Allgemein gilt $H = A + B + C - 2U$ (bei $U$ als Ursprung: $H = A + B + C$): $(8, 4) - (6, 2) = (2, 2)$. \checkmark}
\begin{mitnehmen}\kastentext
Mit $U$ als Ursprung ist $H = A + B + C$ und $S = \frac{A + B + C}3 = \frac H3$ --- die Eulergerade ist dann eine Zeile. Die richtige Wahl des Koordinatenursprungs ersetzt oft eine Seite Rechnung.
\end{mitnehmen}
\end{loesungkarte}

%%%%%%%%%% 6.6
\stupserzwei{tief}{Tangente in $P$ auf Kreis um $M$: $(P - M)\cdot(X - M) = r^{2}$, hier $3(x - 3) + 4(y + 2) = 25$.}{22pt}
\begin{loesungkarte}{tief}{6.6}{Eine Kreisgleichung lesen}
\derweg{Quadratisch ergänzen; Tangente über den Radiusvektor; Schnitt durch Einsetzen; Existenz über das Vorzeichen des Radiusquadrats.}
\schritt{(a)}
$(x^{2} - 6x + 9) + (y^{2} + 4y + 4) = 12 + 9 + 4$: $(x - 3)^{2} + (y + 2)^{2} = 25$. $M = (3, -2)$, $r = 5$.
\schritt{(b)}
$(6 - 3)^{2} + (2 + 2)^{2} = 9 + 16 = 25$. \checkmark\ Tangente: senkrecht zum Radius $P - M = (3, 4)$ durch $P$: $3(x - 6) + 4(y - 2) = 0$, also $3x + 4y = 26$.
\schritt{(c)}
$y = x - 1$ einsetzen: $x^{2} + (x - 1)^{2} - 6x + 4(x - 1) - 12 = 2x^{2} - 4x - 15 = 0$, $x = \frac{4\pm\sqrt{16 + 120}}4 = 1\pm\frac{\sqrt{34}}2$: $x_{1,2} = 3{,}915$ und $-1{,}915$, $y = x - 1$: Schnittpunkte $(3{,}915;\ 2{,}915)$ und $(-1{,}915;\ -2{,}915)$.
\schritt{(d)}
$(x - 3)^{2} + (y + 2)^{2} = 13 - c$. Für $c > 13$ ist die rechte Seite negativ: keine Lösung. Für $c = 13$ nur der Punkt $M$, für $c < 13$ ein Kreis mit $r = \sqrt{13 - c}$.
\probe{(c) Abstand von $M$ zur Geraden $x - y - 1 = 0$: $\frac{|3 + 2 - 1|}{\sqrt2} = 2{,}83 < 5$: zwei Schnittpunkte. \checkmark\ Sehnenlänge $\sqrt2\cdot5{,}83 = 8{,}25$; und $2\sqrt{25 - 8} = 8{,}25$. \checkmark}
\begin{mitnehmen}\kastentext
Kreis: $|X - M| = r$. Tangente in $P$: $(P - M)\cdot(X - M) = r^{2}$. Lage einer Geraden: Abstand von $M$ mit $r$ vergleichen. Alles folgt aus dem Skalarprodukt --- und dasselbe gilt für die Kugel im Raum.
\end{mitnehmen}
\end{loesungkarte}

%%%%%%%%%% 6.7
\stupserzwei{gruen}{$PB = \frac{144}8 = 18$; die Tangentenlänge ist $\sqrt{144} = 12$ --- Pythagoras im Dreieck $P$, Mittelpunkt, Berührpunkt: $13^{2} = 5^{2} + 12^{2}$.}{22pt}
\begin{loesungkarte}{gruen}{6.7}{Die Potenz eines Punktes}
\derweg{Potenz ausrechnen; Sekantensatz und Tangentensatz sind dieselbe Gleichung.}
\schritt{Potenz}
$\operatorname{pot}(P) = d^{2} - r^{2} = 169 - 25 = 144$.
\schritt{Sekante}
$PA\cdot PB = \operatorname{pot}(P)$: $8\cdot PB = 144$, $PB = 18$. Die Sehne $AB$ hat Länge $10$ --- ein Durchmesser; die Sekante geht durch den Mittelpunkt (Probe: $PA = 13 - 5 = 8$). \checkmark
\schritt{Tangente}
$t^{2} = \operatorname{pot}(P) = 144$, $t = 12$. Pythagoras: Berührpunkt $T$, Dreieck $PTM$ rechtwinklig bei $T$: $PT^{2} = PM^{2} - r^{2} = 169 - 25$. \checkmark
\probe{Eine andere Sekante, etwa mit $PA' = 9$: $PB' = 16$, Sehne $7$. Das Produkt bleibt $144$, die Sehne variiert. \checkmark}
\begin{mitnehmen}\kastentext
Die Potenz ist eine Funktion des Punktes: negativ innen (Sehnensatz, $PA\cdot PB = r^{2} - d^{2}$), null auf dem Kreis, positiv außen (Sekanten-Tangenten-Satz). Der Ort der Punkte gleicher Potenz bezüglich zweier Kreise ist eine Gerade, die Potenzgerade --- der Schlüssel zu vielen Kreiskonstruktionen.
\end{mitnehmen}
\end{loesungkarte}

%%%%%%%%%% 6.8
\stupserzwei{gruen}{$PF_1 + PF_2 = 8{,}2 + 1{,}8 = 10 = 2a$. Tangente: $\frac{4x}{25} + \frac{1{,}8y}9 = 1$, also $4x + 5y = 25$.}{22pt}
\begin{loesungkarte}{gruen}{6.8}{Eine Ellipse vermessen}
\derweg{$c$ aus $a^{2} - b^{2}$; Abstände zu den Brennpunkten addieren; Tangentenformel.}
\schritt{(a)}
$c = \sqrt{25 - 9} = 4$, Brennpunkte $(\pm4, 0)$. $e = \frac ca = 0{,}8$. Leitlinien $x = \pm\frac ae = \pm\frac{a^{2}}c = \pm6{,}25$.
\schritt{(b)}
$P$ auf der Ellipse: $\frac{16}{25} + \frac{3{,}24}9 = 0{,}64 + 0{,}36 = 1$. \checkmark\ $PF_1 = \sqrt{(4 + 4)^{2} + 1{,}8^{2}} = \sqrt{67{,}24} = 8{,}2$, $PF_2 = \sqrt{0 + 3{,}24} = 1{,}8$. Summe $10 = 2a$. \checkmark\ Leitlinieneigenschaft: Abstand zu $F_2$ durch Abstand zur Leitlinie $x = 6{,}25$: $\frac{1{,}8}{2{,}25} = 0{,}8 = e$. \checkmark
\schritt{(c)}
Tangente in $(x_0, y_0)$: $\frac{x_0x}{a^{2}} + \frac{y_0y}{b^{2}} = 1$: $\frac{4x}{25} + \frac{1{,}8y}9 = 1$, mal $25$: $4x + 5y = 25$. Steigung $-\frac45$.
\probe{Mit Ableitung: $\frac{2x}{25} + \frac{2yy'}9 = 0$, $y' = -\frac{9x}{25y} = -\frac{36}{45} = -0{,}8$. \checkmark\ Die Tangente halbiert den Außenwinkel zwischen $PF_1$ und $PF_2$ (Reflexionseigenschaft): Richtungen $(-8, -1{,}8)$ und $(0, -1{,}8)$, Tangentenrichtung $(5, -4)$; Winkel zur Tangente: $\cos = \frac{-40 + 7{,}2}{8{,}2\cdot6{,}4} = -0{,}625$ und $\frac{7{,}2}{1{,}8\cdot6{,}4} = 0{,}625$ --- gleich groß. \checkmark}
\begin{mitnehmen}\kastentext
Drei Definitionen, eine Kurve: Summe der Brennpunktabstände $2a$ (Gärtner), Abstandsverhältnis $e < 1$ zur Leitlinie (Kegelschnitt), Gleichung $\frac{x^{2}}{a^{2}} + \frac{y^{2}}{b^{2}} = 1$ (analytisch). Dieselben drei gibt es für Parabel ($e = 1$) und Hyperbel ($e > 1$).
\end{mitnehmen}
\end{loesungkarte}

%%%%%%%%%% 6.9
\stupserzwei{tief}{$u\times v = (1, 1, 0)\times(0, 1, 2) = (2, -2, 1)$, $|u\times v| = 3$. Verbindungsvektor $(-1, 0, 1)$: Abstand $\frac{|-2 + 0 + 1|}3 = \frac13$.}{22pt}
\begin{loesungkarte}{tief}{6.9}{Windschief}
\derweg{Nicht parallel, Gleichsetzen führt zum Widerspruch; Abstand über das Kreuzprodukt; Lotpunkte über zwei Orthogonalitätsbedingungen.}
\schritt{(a)}
Richtungen $(1, 1, 0)$ und $(0, 1, 2)$ sind keine Vielfachen: nicht parallel. Gleichsetzen: $1 + s = 0$, $s = t$, $0 = 1 + 2t$ --- aus den ersten beiden $s = t = -1$, dann $0 = -1$. Kein Schnittpunkt. Windschief.
\schritt{(b)}
$n = u\times v = (2, -2, 1)$, $|n| = 3$. $d = \frac{|(Q - P)\cdot n|}{|n|} = \frac{|(-1, 0, 1)\cdot(2, -2, 1)|}3 = \frac13$.
\schritt{(c)}
$P(s) = (1 + s, s, 0)$, $Q(t) = (0, t, 1 + 2t)$, $Q - P = (-1 - s, t - s, 1 + 2t)$. Senkrecht zu $u$: $-1 - s + t - s = 0$. Senkrecht zu $v$: $t - s + 2 + 4t = 0$. Aus der ersten $t = 1 + 2s$, in die zweite: $5(1 + 2s) - s + 2 = 0$, $s = -\frac79$, $t = -\frac59$. Lotpunkte $P = \bigl(\frac29, -\frac79, 0\bigr)$, $Q = \bigl(0, -\frac59, -\frac19\bigr)$; $Q - P = \bigl(-\frac29, \frac29, -\frac19\bigr)$, Länge $\frac13$. \checkmark
\schritt{(d)}
Ebene durch $(1, 0, 0)$ mit Normale $n$: $2x - 2y + z = 2$. Sie enthält $g$ (Richtung $u\perp n$) und ist parallel zu $h$ ($v\perp n$); $h$ liegt im Abstand $\frac13$ darüber.
\probe{$Q - P$ ist parallel zu $n$: $\bigl(-\frac29, \frac29, -\frac19\bigr) = -\frac19(2, -2, 1)$. \checkmark}
\begin{mitnehmen}\kastentext
Zwei Geraden im Raum: Richtungen vergleichen (parallel?), Gleichsetzen (Schnitt?), sonst windschief. Der Abstand ist die Projektion des Verbindungsvektors auf die gemeinsame Normale $u\times v$ --- dieselbe Formel wie der Abstand Punkt--Ebene, denn $h$ liegt in einer zu $g$ parallelen Ebene.
\end{mitnehmen}
\end{loesungkarte}

%%%%%%%%%% 6.10
\stupserzwei{akzent}{(a) $54{,}7^{\circ}$ ist der Winkel zur \emph{Normalen}. Der Winkel zur Ebene ist $90^{\circ} - 54{,}7^{\circ} = 35{,}3^{\circ}$, und $\sin35{,}3^{\circ} = \frac1{\sqrt3}$.}{22pt}
\begin{loesungkarte}{akzent}{6.10}{Sinus oder Kosinus?}
\derweg{Jede Formel auf die Frage zurückführen: Winkel zwischen welchen zwei Vektoren? Richtung und Normale sind verschiedene Dinge.}
\schritt{(a)}
Das Skalarprodukt mit der Normalen gibt den Winkel zwischen Gerade und \emph{Normale}: $54{,}7^{\circ}$. Der Winkel zwischen Gerade und Ebene ist das Komplement: $35{,}3^{\circ}$. Formel: $\sin\varphi = \frac{|d\cdot n|}{|d||n|}$ --- Sinus, weil $\varphi = 90^{\circ} - \angle(d, n)$.
\schritt{(b)}
Winkel zwischen Ebenen $=$ Winkel zwischen Normalen: $\cos\varphi = \frac{|n_1\cdot n_2|}{|n_1||n_2|} = \frac{|(1, 1, 0)\cdot(0, 0, 1)|}{\sqrt2} = 0$, $\varphi = 90^{\circ}$. Zwei Ebenen bilden zwei Winkel ($\varphi$ und $180^{\circ} - \varphi$); gemeint ist der spitze --- daher der Betrag.
\schritt{(c)}
$\cos\varphi = \frac{(1, 0, 0)\cdot(1, 1, 0)}{1\cdot\sqrt2} = \frac1{\sqrt2}$, $\varphi = 45^{\circ}$. Mit $h'$: $\frac{-1}{\sqrt2}$, $135^{\circ}$ --- dieselbe Gerade, anderer Richtungsvektor. Der Winkel zwischen Geraden ist in $[0^{\circ}, 90^{\circ}]$: Betrag setzen, $45^{\circ}$.
\schritt{(d)}
\begin{center}\small
\begin{tabular}{l|l|l}
Objekte & Vektoren & Formel\\ \hline
Gerade--Gerade & Richtungen $d_1, d_2$ & $\cos\varphi = \frac{|d_1\cdot d_2|}{|d_1||d_2|}$\\
Ebene--Ebene & Normalen $n_1, n_2$ & $\cos\varphi = \frac{|n_1\cdot n_2|}{|n_1||n_2|}$\\
Gerade--Ebene & Richtung $d$, Normale $n$ & $\sin\varphi = \frac{|d\cdot n|}{|d||n|}$
\end{tabular}
\end{center}
\begin{fehler}\kastentext
Das Skalarprodukt misst den Winkel zwischen den \emph{Vektoren}, die man einsetzt. Bei Gerade und Ebene setzt man eine Richtung und eine Normale ein --- und bekommt das Komplement des gesuchten Winkels. Der Betrag im Zähler macht aus zwei möglichen Winkeln den spitzen.
\end{fehler}
\begin{mitnehmen}\kastentext
Vor jeder Winkelrechnung: Welche zwei Vektoren beschreiben die Objekte, und ist einer davon eine Normale? Eine Normale im Spiel: Vorzeichen des Komplements beachten. Zwei Normalen: Kosinus wie bei Geraden.
\end{mitnehmen}
\end{loesungkarte}

%%%%%%%%%% 6.11
\stupserzwei{gruen}{$12x + 7y + 11z = 59$; $D$ eingesetzt gibt $51$: Abstand $\frac8{\sqrt{314}} = 0{,}451$. Volumen $\frac86 = \frac43$.}{22pt}
\begin{loesungkarte}{gruen}{6.11}{Ebene, Abstand, Tetraeder}
\derweg{Normale aus dem Kreuzprodukt; Abstand als normierte Projektion; Volumen als Spatprodukt durch $6$; Winkel über das Skalarprodukt.}
\schritt{(a)}
$u = B - A = (3, -2, -2)$, $v = C - A = (1, 3, -3)$. $n = u\times v = \bigl((-2)(-3) - (-2)(3),\ (-2)(1) - (3)(-3),\ 3\cdot3 - (-2)(1)\bigr) = (12, 7, 11)$. Ebene: $12x + 7y + 11z = 12 + 14 + 33 = 59$. Probe mit $B$: $48 + 0 + 11 = 59$ \checkmark, $C$: $24 + 35 + 0 = 59$ \checkmark.
\schritt{(b)}
$D$ eingesetzt: $0 + 7 + 44 = 51$. $|n| = \sqrt{144 + 49 + 121} = \sqrt{314} = 17{,}72$. Abstand $\frac{|51 - 59|}{17{,}72} = 0{,}451$.
\schritt{(c)}
$V = \frac16|n\cdot(D - A)| = \frac16|(12, 7, 11)\cdot(-1, -1, 1)| = \frac16|-12 - 7 + 11| = \frac86 = \frac43$.
\schritt{(d)}
$\cos\varphi = \frac{u\cdot v}{|u||v|} = \frac{3 - 6 + 6}{\sqrt{17}\sqrt{19}} = \frac3{17{,}97} = 0{,}167$, $\varphi = 80{,}4^{\circ}$.
\probe{(c) anders: Grundfläche $\frac12|n| = 8{,}86$, Höhe $0{,}451$: $V = \frac13\cdot8{,}86\cdot0{,}451 = 1{,}33$. \checkmark}
\begin{mitnehmen}\kastentext
Kreuzprodukt: Normale und Flächeninhalt ($|u\times v|$ = Parallelogramm). Spatprodukt: Volumen ($|n\cdot w|$ = Spat, durch $6$ = Tetraeder). Skalarprodukt: Winkel und Abstände. Drei Produkte, und die ganze Raumgeometrie der Schule ist damit erledigt.
\end{mitnehmen}
\end{loesungkarte}

%%%%%%%%%% 6.12
\stupserzwei{akzent}{Die erweiterte $3\times3$-Matrix hat bei echten Kegelschnitten Determinante $\neq0$; bei (a), (c), (d), (e) ist sie $0$. (b) ist nicht entartet, nur leer.}{22pt}
\begin{loesungkarte}{akzent}{6.12}{Entartete Kegelschnitte}
\derweg{Faktorisieren, wo es geht; sonst Vorzeichen betrachten; die Determinante der erweiterten Matrix als Kriterium.}
\schritt{Die fünf Mengen}
(a) $(x - y)(x + y) = 0$: zwei sich schneidende Geraden $y = \pm x$. \\
(b) $x^{2} + y^{2} = -1$: leer (über $\mathbb R$). \\
(c) $(x - y)^{2} = 0$: eine Gerade $y = x$, doppelt gezählt. \\
(d) $(x - y)^{2} = 1$: zwei parallele Geraden $y = x\pm1$. \\
(e) $xy = 0$: die beiden Koordinatenachsen.
\schritt{Echt oder entartet}
Echte Kegelschnitte: Ellipse, Parabel, Hyperbel. (a), (c), (d), (e) sind Geradenpaare --- entartet. (b) ist keine Kurve; man nennt sie \enquote{imaginäre Ellipse}: algebraisch nicht entartet (die Gleichung ist irreduzibel), aber ohne reelle Punkte.
\schritt{Kriterium}
Schreibe $x^{\mathsf T}Ax + 2b^{\mathsf T}x + c = 0$ mit der erweiterten Matrix $\tilde A = \begin{pmatrix}A & b\\ b^{\mathsf T} & c\end{pmatrix}$. Entartet $\iff\det\tilde A = 0$. (a) $\tilde A = \operatorname{diag}(1, -1, 0)$: $0$. (c) $A = \begin{pmatrix}1 & -1\\ -1 & 1\end{pmatrix}$, $c = 0$: $\det\tilde A = 0$. (d) dasselbe $A$, $c = -1$: $\det = -1\cdot\det A = 0$. (e) $A = \begin{pmatrix}0 & \frac12\\ \frac12 & 0\end{pmatrix}$, $c = 0$: $0$. (b) $\operatorname{diag}(1, 1, 1)$: $1\neq0$, nicht entartet. Die kleine Determinante $\det A$ sagt nur den Typ: $> 0$ elliptisch, $< 0$ hyperbolisch, $= 0$ parabolisch --- (c) und (d) sind \enquote{parabolisch entartet}.
\probe{Eine echte Hyperbel $x^{2} - y^{2} = 1$: $\tilde A = \operatorname{diag}(1, -1, -1)$, $\det = 1\neq0$. \checkmark\ Ihre Asymptoten sind genau der entartete Fall (a). \checkmark}
\begin{fehler}\kastentext
Vor der Hauptachsentransformation: $\det\tilde A$ prüfen. Ist sie null, zerfällt die Gleichung in lineare Faktoren (oder einen Punkt), und Halbachsen gibt es nicht. Und eine Gleichung ohne reelle Lösung ist kein Widerspruch --- $x^{2} + y^{2} = -1$ ist eine korrekte Gleichung mit leerer Lösungsmenge.
\end{fehler}
\begin{mitnehmen}\kastentext
Entartete Kegelschnitte sind die Grenzfälle der echten: Zwei sich schneidende Geraden sind eine Hyperbel mit Halbachsen $0$, zwei parallele eine Parabel ohne Scheitel, ein Punkt eine Ellipse mit Halbachsen $0$. In der projektiven Geometrie (Kapitel 7) sind sie alle Schnitte eines Kegels mit einer Ebene --- durch die Spitze.
\end{mitnehmen}
\end{loesungkarte}

%%%%%%%%%% 6.13
\stupserzwei{gruen}{Eigenvektoren: zu $4$ ist $(2, 1)$, zu $9$ ist $(1, -2)$. In diesen Richtungen liegen die Halbachsen $3$ und $2$.}{22pt}
\begin{loesungkarte}{gruen}{6.13}{Typ und Normalform}
\derweg{Matrix, Determinante für den Typ, Eigenwerte und -vektoren, Normalform.}
\schritt{Matrix}
$5x^{2} - 4xy + 8y^{2} = (x, y)\begin{pmatrix}5 & -2\\ -2 & 8\end{pmatrix}\begin{pmatrix}x\\ y\end{pmatrix}$.
\schritt{Typ}
$\det A = 40 - 4 = 36 > 0$: elliptisch. Erweiterte Determinante $\det\tilde A = -36\cdot36\neq0$: nicht entartet. Da $5 > 0$ und rechte Seite $36 > 0$: eine echte Ellipse.
\schritt{Eigenwerte}
$\lambda^{2} - 13\lambda + 36 = 0$: $\lambda_1 = 4$, $\lambda_2 = 9$. Eigenvektoren: $(A - 4I)v = 0$ mit $\begin{pmatrix}1 & -2\\ -2 & 4\end{pmatrix}$: $v_1 = \frac1{\sqrt5}(2, 1)$. $(A - 9I)$: $\begin{pmatrix}-4 & -2\\ -2 & -1\end{pmatrix}$: $v_2 = \frac1{\sqrt5}(1, -2)$.
\schritt{Normalform}
$4x'^{2} + 9y'^{2} = 36$, also $\frac{x'^{2}}9 + \frac{y'^{2}}4 = 1$: Ellipse mit Halbachsen $a = 3$ in Richtung $(2, 1)$ und $b = 2$ in Richtung $(1, -2)$. Drehwinkel $\arctan\frac12 = 26{,}6^{\circ}$.
\probe{Scheitel $3v_1 = \frac3{\sqrt5}(2, 1) = (2{,}68;\ 1{,}34)$: $5\cdot7{,}2 - 4\cdot3{,}6 + 8\cdot1{,}8 = 36 - 14{,}4 + 14{,}4 = 36$. \checkmark\ Spur $13 = 4 + 9$, Determinante $36 = 4\cdot9$. \checkmark}
\begin{mitnehmen}\kastentext
Die Eigenwerte sagen den Typ (beide positiv: Ellipse; verschiedene Vorzeichen: Hyperbel; einer null: Parabel), die Eigenvektoren die Achsen, die Quotienten $\sqrt{c/\lambda_i}$ die Halbachsen. Der große Eigenwert gehört zur kurzen Achse.
\end{mitnehmen}
\end{loesungkarte}

%%%%%%%%%% 6.14
\stupserzwei{gruen}{$4x'^{2} + y'^{2} + z'^{2} = 6$: Ellipsoid mit Halbachsen $\sqrt{1{,}5}$, $\sqrt6$, $\sqrt6$ --- ein abgeplattetes Rotationsellipsoid um die Achse $(1, 1, 1)$.}{22pt}
\begin{loesungkarte}{gruen}{6.14}{Eine Quadrik im Raum}
\derweg{Struktur der Matrix erkennen, Eigenwerte ablesen, Normalform, Halbachsen.}
\schritt{Matrix}
$A = \begin{pmatrix}2 & 1 & 1\\ 1 & 2 & 1\\ 1 & 1 & 2\end{pmatrix} = I + J$ mit $J$ der Matrix aus lauter Einsen.
\schritt{Eigenwerte}
$J$ hat Rang $1$: Eigenwert $3$ zum Vektor $(1, 1, 1)$ und Eigenwert $0$ (doppelt) auf dem Komplement. Also hat $A = I + J$ die Eigenwerte $4$ (einfach, Richtung $(1, 1, 1)$) und $1$ (doppelt, Ebene $x + y + z = 0$).
\schritt{Normalform}
$4x'^{2} + y'^{2} + z'^{2} = 6$, also $\frac{x'^{2}}{3/2} + \frac{y'^{2}}6 + \frac{z'^{2}}6 = 1$. Alle Eigenwerte positiv: Ellipsoid. Halbachsen $\sqrt{1{,}5} = 1{,}22$ in Richtung $(1, 1, 1)$, $\sqrt6 = 2{,}45$ in jeder Richtung senkrecht dazu. Ein abgeplattetes Rotationsellipsoid (wie die Erde), Symmetrieachse $(1, 1, 1)$.
\probe{Punkt $\frac{1{,}22}{\sqrt3}(1, 1, 1) = (0{,}707, 0{,}707, 0{,}707)$: $3\cdot2\cdot0{,}5 + 3\cdot2\cdot0{,}5 = 3 + 3 = 6$. \checkmark\ Punkt $\sqrt3\,(1, -1, 0)$ (Länge $\sqrt6$): $2\cdot3 + 2\cdot3 + 0 + 2(-3) + 0 + 0 = 6$. \checkmark}
\begin{mitnehmen}\kastentext
Matrizen der Form $aI + bJ$ haben Eigenwerte $a + nb$ (einmal) und $a$ ($n - 1$-mal) --- das sieht man ohne Rechnung, und solche Matrizen sind überall: Kovarianzen gleich korrelierter Größen, vollständige Graphen, Austauschsymmetrie. Doppelte Eigenwerte heißen Rotationssymmetrie.
\end{mitnehmen}
\end{loesungkarte}

%%%%%%%%%% 6.15
\stupserzwei{tief}{$f_2$: Determinante $-1$, keine Fixpunkte --- Gleitspiegelung. Die Achse geht durch die Mittelpunkte von $P$ und $f_2(P)$: $(0, 0)\mapsto(1, 1)$, Mittelpunkt $(\frac12, \frac12)$.}{22pt}
\begin{loesungkarte}{tief}{6.15}{Drei Abbildungen, drei Typen}
\derweg{Linearen Teil ablesen, Determinante, Fixpunktgleichung lösen; was übrig bleibt, bestimmt den Typ.}
\schritt{$f_1$: Drehung}
Linearer Teil $\begin{pmatrix}0 & -1\\ 1 & 0\end{pmatrix}$, Determinante $+1$, Drehmatrix um $90^{\circ}$. Fixpunkt: $x = -y + 2$, $y = x + 1$, also $x = -(x + 1) + 2$, $x = \frac12$, $y = \frac32$. Drehung um $90^{\circ}$ gegen den Uhrzeigersinn um $Z = \bigl(\frac12, \frac32\bigr)$.
\schritt{$f_2$: Gleitspiegelung}
Linearer Teil $\begin{pmatrix}0 & 1\\ 1 & 0\end{pmatrix}$, Determinante $-1$: Spiegelungsmatrix an $y = x$. Fixpunkte: $x = y + 1$ und $y = x + 1$ --- Widerspruch, keine. Also Gleitspiegelung. Achse: Mittelpunkte von $P$ und $f_2(P)$ liegen darauf: $(0, 0)\mapsto(1, 1)$, Mitte $(\frac12, \frac12)$; $(1, 0)\mapsto(1, 2)$, Mitte $(1, 1)$. Achse $y = x$. Gleitvektor: $f_2$ verschiebt Punkte der Achse um $(1, 1)$.
\schritt{$f_3$: Spiegelung}
Linearer Teil $\begin{pmatrix}-1 & 0\\ 0 & 1\end{pmatrix}$, Determinante $-1$. Fixpunkte: $x = -x + 4$, also $x = 2$, $y$ beliebig: eine ganze Gerade. Spiegelung an $x = 2$.
\probe{$f_1$ viermal angewandt auf $(0, 0)$: $(2, 1)\to(1, 3)\to(-1, 2)\to(0, 0)$. \checkmark\ $f_2$ zweimal: $(0, 0)\to(1, 1)\to(2, 2)$ --- Verschiebung um $2\cdot(1, 1)$, wie bei einer Gleitspiegelung. \checkmark\ $f_3$ zweimal: Identität. \checkmark}
\begin{mitnehmen}\kastentext
Entscheidungsbaum: $\det = +1$ und Fixpunkt $\Rightarrow$ Drehung; $\det = +1$ ohne Fixpunkt $\Rightarrow$ Verschiebung; $\det = -1$ mit Fixgerade $\Rightarrow$ Spiegelung; $\det = -1$ ohne Fixpunkt $\Rightarrow$ Gleitspiegelung. Vier Typen, und jede Isometrie der Ebene ist einer davon (Chasles).
\end{mitnehmen}
\end{loesungkarte}

%%%%%%%%%% 6.16
\stupserzwei{gruen}{$S_2S_1 = \begin{pmatrix}0 & 1\\ 1 & 0\end{pmatrix}\begin{pmatrix}1 & 0\\ 0 & -1\end{pmatrix} = \begin{pmatrix}0 & -1\\ 1 & 0\end{pmatrix}$: Drehung um $+90^{\circ}$. Umgekehrt: $-90^{\circ}$.}{22pt}
\begin{loesungkarte}{gruen}{6.16}{Zwei Spiegelungen}
\derweg{Matrizen aufstellen, in der richtigen Reihenfolge multiplizieren, Ergebnis als Drehmatrix erkennen.}
\schritt{Matrizen}
Spiegelung an $y = 0$: $S_1 = \begin{pmatrix}1 & 0\\ 0 & -1\end{pmatrix}$. Spiegelung an $y = x$: $S_2 = \begin{pmatrix}0 & 1\\ 1 & 0\end{pmatrix}$.
\schritt{Verkettung}
Erst $S_1$, dann $S_2$: $S_2S_1 = \begin{pmatrix}0 & -1\\ 1 & 0\end{pmatrix} = \begin{pmatrix}\cos90^{\circ} & -\sin90^{\circ}\\ \sin90^{\circ} & \cos90^{\circ}\end{pmatrix}$: Drehung um $90^{\circ}$ gegen den Uhrzeigersinn um den Ursprung (den Schnittpunkt der Achsen). Allgemein: Winkel zwischen den Achsen $45^{\circ}$, Drehwinkel das Doppelte.
\schritt{Umgekehrt}
$S_1S_2 = \begin{pmatrix}0 & 1\\ -1 & 0\end{pmatrix}$: Drehung um $-90^{\circ}$. Die Reihenfolge der Spiegelungen bestimmt den Drehsinn; Spiegelungen kommutieren nur, wenn die Achsen senkrecht stehen (dann ist das Produkt die Punktspiegelung).
\probe{$(1, 0)\xrightarrow{S_1}(1, 0)\xrightarrow{S_2}(0, 1)$: um $90^{\circ}$ gedreht. \checkmark\ $(S_2S_1)^{4} = I$. \checkmark}
\begin{mitnehmen}\kastentext
Zwei Spiegelungen an Geraden mit Schnittwinkel $\varphi$ sind eine Drehung um $2\varphi$ um den Schnittpunkt; an parallelen Geraden im Abstand $d$ eine Verschiebung um $2d$. Deshalb erzeugen Spiegelungen alle Isometrien --- und jede Isometrie der Ebene ist Produkt von höchstens drei.
\end{mitnehmen}
\end{loesungkarte}

%%%%%%%%%% 6.17
\stupserzwei{gruen}{Weg 2: Die vier Raumdiagonalen werden von jeder Drehung permutiert, und jede der $24$ Permutationen kommt vor. Mit der Punktspiegelung ($-I$) verdoppelt sich alles.}{22pt}
\begin{loesungkarte}{gruen}{6.17}{Die Symmetriegruppe des Würfels}
\derweg{Bahnformel für die Drehgruppe; Permutationen der Raumdiagonalen als zweiter Weg; dann Spiegelungen dazu.}
\schritt{Weg 1: Bahnformel}
Die Drehgruppe operiert auf den $6$ Seitenflächen. Bahn einer Fläche: alle $6$ (jede lässt sich in jede drehen). Stabilisator: die $4$ Drehungen um die Achse durch den Flächenmittelpunkt ($0^{\circ}, 90^{\circ}, 180^{\circ}, 270^{\circ}$). $|G_{\text{Dreh}}| = 6\cdot4 = 24$.
\schritt{Weg 2: Raumdiagonalen}
Jede Drehung permutiert die $4$ Raumdiagonalen; verschiedene Drehungen geben verschiedene Permutationen (nur die Identität lässt alle vier fest). Das liefert eine injektive Abbildung $G_{\text{Dreh}}\to S_4$, und $|S_4| = 24$. Da man alle Transpositionen durch Drehungen um $180^{\circ}$ um Kantenmittenachsen bekommt, ist sie surjektiv: $G_{\text{Dreh}}\cong S_4$.
\schritt{Volle Gruppe}
Die Punktspiegelung $-I$ ist eine Symmetrie des Würfels, keine Drehung, und vertauscht mit allen Drehungen. Jede Symmetrie ist eine Drehung oder eine Drehung mal $-I$: $|G| = 2\cdot24 = 48$, $G\cong S_4\times\mathbb Z_2$.
\schritt{Zählen nach Typ}
Drehungen: $1$ Identität; $9$ um Flächenachsen ($3$ Achsen, je $3$ Winkel); $8$ um Raumdiagonalen ($4$ Achsen, je $2$ Winkel $\pm120^{\circ}$); $6$ um Kantenmittenachsen ($180^{\circ}$). $1 + 9 + 8 + 6 = 24$. \checkmark
\probe{Bahnformel mit Ecken: Bahn $8$, Stabilisator $3$ (Drehungen um die Raumdiagonale durch die Ecke): $24$. \checkmark\ Mit Kanten: $12\cdot2 = 24$. \checkmark}
\begin{mitnehmen}\kastentext
Die Bahnformel $|G| = |\text{Bahn}|\cdot|\text{Stabilisator}|$ zählt Gruppen, ohne ihre Elemente aufzulisten; Oktaeder (dual) hat dieselbe Gruppe, Tetraeder $S_4$ als volle Gruppe mit $12$ Drehungen, Ikosaeder $A_5\times\mathbb Z_2$ mit $120$ Elementen. Diese Gruppen sind die Symmetrien der Moleküle und Kristalle (Band 30).
\end{mitnehmen}
\end{loesungkarte}

%%%%%%%%%% 6.18
\stupserzwei{tief}{$t = (1, 1)$; $M e_1 = (3, 2) - (1, 1) = (2, 1)$, $M e_2 = (2, 4) - (1, 1) = (1, 3)$. $\det M = 5$: Flächen werden verfünffacht.}{22pt}
\begin{loesungkarte}{tief}{6.18}{Eine affine Abbildung bestimmen}
\derweg{Verschiebung aus dem Bild des Ursprungs, Matrixspalten aus den Bildern der Einheitsvektoren; dann Invarianten prüfen.}
\schritt{(a)}
$t = f(0, 0) = (1, 1)$. $M = \begin{pmatrix}2 & 1\\ 1 & 3\end{pmatrix}$. $f(x, y) = (2x + y + 1,\ x + 3y + 1)$.
\schritt{(b)}
Ecken $(0, 0), (1, 0), (1, 1), (0, 1)\mapsto(1, 1), (3, 2), (4, 5), (2, 4)$: ein Parallelogramm. Fläche $|\det M|\cdot1 = 5$. Probe: Vektoren $(2, 1)$ und $(1, 3)$, Determinante $6 - 1 = 5$. \checkmark
\schritt{(c)}
Mittelpunkt von $(0, 0)$ und $(1, 0)$ ist $(\frac12, 0)$, Bild $(2, \frac32)$ --- der Mittelpunkt von $(1, 1)$ und $(3, 2)$. \checkmark\ Die Geraden $y = 0$ und $y = 1$ werden zu $(1, 1) + s(2, 1)$ und $(2, 4) + s(2, 1)$: parallel. \checkmark\ Der rechte Winkel bei $(0, 0)$ wird zum Winkel zwischen $(2, 1)$ und $(1, 3)$: $\cos = \frac5{\sqrt5\sqrt{10}} = 0{,}707$, also $45^{\circ}$. Nicht erhalten.
\schritt{(d)}
$f^{-1}(x) = M^{-1}(x - t)$ mit $M^{-1} = \frac15\begin{pmatrix}3 & -1\\ -1 & 2\end{pmatrix}$ --- wieder affin. Affine Abbildungen bilden eine Gruppe.
\probe{$f^{-1}(4, 5) = \frac15\begin{pmatrix}3 & -1\\ -1 & 2\end{pmatrix}\begin{pmatrix}3\\ 4\end{pmatrix} = \frac15(5, 5) = (1, 1)$. \checkmark\ Das Parallelogramm: $(3, 2) - (1, 1) = (2, 1)$ und $(4, 5) - (2, 4) = (2, 1)$, gegenüberliegende Seiten gleich. \checkmark}
\begin{mitnehmen}\kastentext
Affin erhalten: Geraden, Parallelität, Teilverhältnisse, Flächen\emph{verhältnisse} (Faktor $|\det M|$ für alle). Nicht erhalten: Längen, Winkel, Kreise (werden Ellipsen). Was erhalten bleibt, darf man am bequemsten Dreieck beweisen (Aufgabe 6.20).
\end{mitnehmen}
\end{loesungkarte}

%%%%%%%%%% 6.19
\stupserzwei{gruen}{$\frac{AC}{CB} = \frac24 = \frac12$; also muss $\frac{AD}{DB} = -\frac12$: $\frac x{6 - x} = -\frac12$, $x = -6$.}{22pt}
\begin{loesungkarte}{gruen}{6.19}{Harmonische Lage}
\derweg{Doppelverhältnis mit orientierten Längen hinschreiben und gleich $-1$ setzen.}
\schritt{Rechnung}
$(A, B;\ C, D) = \frac{AC}{CB}\Big/\frac{AD}{DB} = \frac{2 - 0}{6 - 2}\Big/\frac{x - 0}{6 - x} = \frac12\cdot\frac{6 - x}x$. Gleich $-1$: $6 - x = -2x$, $x = -6$. $D = -6$.
\schritt{Deutung}
$C = 2$ teilt $AB$ innen im Verhältnis $AC : CB = 2 : 4 = 1 : 2$. $D = -6$ teilt außen: $|AD| : |DB| = 6 : 12 = 1 : 2$. Gleiches Verhältnis, einmal innen, einmal außen --- das ist harmonische Teilung.
\probe{Vertauscht man die Paare: $(C, D;\ A, B) = \frac{CA}{AD}\Big/\frac{CB}{BD} = \frac{-2}{-6}\Big/\frac4{-12} = \frac13\Big/\bigl(-\frac13\bigr) = -1$. \checkmark\ Harmonische Lage ist symmetrisch in den beiden Paaren.}
\begin{mitnehmen}\kastentext
Das Doppelverhältnis ist die einzige projektive Invariante von vier Punkten auf einer Geraden; $-1$ ist der ausgezeichnete Wert, bei dem die Vierergruppe symmetrisch wird. Harmonische Punkte entstehen beim vollständigen Vierseit, bei Pol und Polare am Kreis --- und beim Mittelpunkt: $M$ und der Fernpunkt teilen $AB$ harmonisch.
\end{mitnehmen}
\end{loesungkarte}

%%%%%%%%%% 6.20
\stupserzwei{gruen}{Mit $A = (0, 0)$, $B = e_1$, $C = e_2$ ist $U = I$: $M = U'$ hat die Spalten $B' - A'$ und $C' - A'$, und $t = A'$.}{22pt}
\begin{loesungkarte}{gruen}{6.20}{Affin gibt es nur ein Dreieck}
\derweg{Kantenvektoren als Basen; $M$ überführt die eine Basis in die andere; $t$ aus $A\mapsto A'$.}
\schritt{Allgemein}
$u = B - A$, $v = C - A$; $u', v'$ ebenso. Nichtentartet heißt: $\{u, v\}$ und $\{u', v'\}$ sind Basen. Sei $U$ die Matrix mit Spalten $u, v$ und $U'$ die mit $u', v'$; beide invertierbar. $M = U'U^{-1}$ erfüllt $Mu = u'$, $Mv = v'$. Mit $t = A' - MA$ ist $f(x) = Mx + t$ affin und $f(A) = A'$, $f(B) = MA + Mu + t = A' + u' = B'$, $f(C) = C'$. $\blacksquare$ Eindeutig, weil eine affine Abbildung durch drei nichtkollineare Punkte festgelegt ist.
\schritt{Konkret}
$U = I$, also $M = U' = \begin{pmatrix}3 & 1\\ 1 & 4\end{pmatrix}$ (Spalten $B' - A' = (3, 1)$, $C' - A' = (1, 4)$), $t = A' = (1, 1)$. $f(x, y) = (3x + y + 1,\ x + 4y + 1)$.
\schritt{Schwerpunkt}
$S = \bigl(\frac13, \frac13\bigr)$, $f(S) = \bigl(1 + \frac13 + 1,\ \frac13 + \frac43 + 1\bigr) = \bigl(\frac73, \frac83\bigr) = \frac{A' + B' + C'}3$. Affine Abbildungen erhalten Schwerpunkte --- weil der Schwerpunkt über Teilverhältnisse definiert ist. Deshalb genügt es, \enquote{die Seitenhalbierenden schneiden sich in einem Punkt und teilen sich $2 : 1$} am gleichseitigen Dreieck zu beweisen.
\probe{$f(1, 0) = (4, 2) = B'$, $f(0, 1) = (2, 5) = C'$. \checkmark\ $\det M = 11 > 0$: orientierungserhaltend; die Fläche von $A'B'C'$ ist $\frac{11}2$, die von $ABC$ ist $\frac12$. \checkmark}
\begin{mitnehmen}\kastentext
Affine Geometrie kennt kein \enquote{gleichseitig} und kein \enquote{rechtwinklig} --- alle Dreiecke sind äquivalent, alle Parallelogramme, alle Ellipsen. Sätze über Seitenhalbierende, Mittelparallelen, Schwerpunkte, Teilverhältnisse sind affin und dürfen am Spezialfall bewiesen werden. Sätze über Höhen, Umkreis, Winkel sind es nicht.
\end{mitnehmen}
\end{loesungkarte}

%%%%%%%%%% 6.21
\stupserzwei{akzent}{(c) Äquivalent zum Parallelenaxiom: Rechteck, ähnliche nichtkongruente Dreiecke, Pythagoras, Umkreis für jedes Dreieck. In jeder Geometrie: Winkelsumme $\le180^{\circ}$ (Saccheri--Legendre).}{22pt}
\begin{loesungkarte}{akzent}{6.21}{Ein Beweis im Kreis}
\derweg{Prüfen, welche Voraussetzung an welcher Stelle eingeht; die Äquivalenzen auseinanderhalten; das Modell als Unabhängigkeitsbeweis.}
\schritt{(a)}
Die Winkelsumme $180^{\circ}$ ist dem Parallelenaxiom äquivalent: Man kann sie nur \emph{mit} dem Axiom beweisen. Wer sie voraussetzt, hat das Axiom schon benutzt --- der Beweis ist zirkulär. (In der hyperbolischen Ebene ist die Winkelsumme kleiner, und durch $P$ gibt es viele Parallelen; beides zusammen ist konsistent.)
\schritt{(b)}
Dreieck $ABC$; durch $C$ die Parallele $p$ zu $AB$ (existiert in jeder Geometrie; \emph{eindeutig} nach dem Axiom). Wechselwinkel an den Transversalen $AC$ und $BC$: Die Winkel bei $C$ zwischen $p$ und $CA$ bzw.\ $CB$ sind $\alpha$ und $\beta$ --- das braucht die Eindeutigkeit: Nur die eine Parallele hat gleiche Wechselwinkel. Zusammen mit $\gamma$ füllen sie den gestreckten Winkel bei $C$: $\alpha + \beta + \gamma = 180^{\circ}$. $\blacksquare$
\schritt{(c)}
Äquivalent zum Parallelenaxiom: Existenz eines Rechtecks (vier rechte Winkel); Existenz ähnlicher, nicht kongruenter Dreiecke (Wallis); Satz des Pythagoras; Umkreis für jedes Dreieck (Bolyai). Unabhängig davon gültig (absolute Geometrie): Winkelsumme $\le180^{\circ}$ (Saccheri--Legendre), Außenwinkelsatz (Außenwinkel $>$ jeder nicht anliegende Innenwinkel), Existenz \emph{mindestens einer} Parallelen.
\schritt{(d)}
Die hyperbolische Ebene (Poincaré-Modell) erfüllt alle Axiome Euklids außer dem Parallelenaxiom. Wäre das Axiom aus den anderen beweisbar, müsste es auch dort gelten --- tut es nicht. Also ist es unabhängig. Das ist ein Beweis über ein Modell: Man zeigt nicht, dass etwas wahr ist, sondern dass seine Negation konsistent ist.
\begin{fehler}\kastentext
Ein Zirkelschluss versteckt sich meist in einer \enquote{offensichtlichen} Aussage, die der Behauptung äquivalent ist. Test: Gilt die benutzte Aussage auch in der hyperbolischen Ebene? Wenn nein, hat man das Parallelenaxiom benutzt --- so scheiterten Proklos, Saccheri, Lambert und Legendre.
\end{fehler}
\begin{mitnehmen}\kastentext
Unabhängigkeit beweist man durch Modelle: ein Modell mit dem Axiom (euklidische Ebene), eines ohne (hyperbolische Ebene). Dasselbe Verfahren entscheidet in der Mengenlehre über das Auswahlaxiom und die Kontinuumshypothese (Band 10).
\end{mitnehmen}
\end{loesungkarte}

%%%%%%%%%% 6.22
\stupserzwei{tief}{(d) Winkelsumme $180^{\circ}$ gäbe Fläche $0$, Winkelsumme $210^{\circ}$ Fläche negativ: beide unmöglich. Gleichseitige hyperbolische Dreiecke haben Winkel $< 60^{\circ}$.}{22pt}
\begin{loesungkarte}{tief}{6.22}{Hyperbolische Vielecke}
\derweg{Fläche $=$ Defekt: $(n - 2)\pi$ minus Winkelsumme; alles andere ist Auflösen.}
\schritt{(a)}
$\alpha + \beta + \gamma = \pi - 1$. $\gamma = \pi - 1 - \frac{50\pi}{180} = 3{,}1416 - 1 - 0{,}8727 = 1{,}2689$ rad $= 72{,}7^{\circ}$. Winkelsumme $122{,}7^{\circ}$.
\schritt{(b)}
Fünfeck: $(5 - 2)\pi - 5\cdot\frac\pi2 = 3\pi - 2{,}5\pi = \frac\pi2$. Euklidisch wäre die Winkelsumme eines Fünfecks $540^{\circ}$, die eines rechtwinkligen aber $450^{\circ}$ --- Widerspruch; hyperbolisch ist das Defizit von $90^{\circ}$ genau die Fläche.
\schritt{(c)}
Fläche $= \pi - (\alpha + \beta + \gamma) < \pi$. Der Grenzwert $\pi$ wird vom idealen Dreieck erreicht: drei Ecken im Unendlichen, drei Winkel $0$, unendlich lange Seiten --- und endliche Fläche $\pi$.
\schritt{(d)}
$60^{\circ}, 60^{\circ}, 60^{\circ}$: Summe $\pi$, Fläche $0$ --- nur als entartetes Dreieck. $70^{\circ}$: Summe $> \pi$, Fläche negativ --- unmöglich. Gleichseitige hyperbolische Dreiecke haben Winkel echt kleiner als $60^{\circ}$, umso kleiner, je größer sie sind. Es gibt kein \emph{das} gleichseitige Dreieck: Größe und Winkel hängen zusammen (keine Ähnlichkeit, Aufgabe 6.21).
\probe{(b) Fünf rechtwinklige Fünfecke um eine Ecke? Nein --- $4\cdot90^{\circ} = 360^{\circ}$: vier. Die hyperbolische Ebene lässt sich mit rechtwinkligen Fünfecken parkettieren, je vier an jeder Ecke. \checkmark}
\begin{mitnehmen}\kastentext
Hyperbolisch: Fläche $= (n - 2)\pi - \text{Winkelsumme}$ (bei $K = -1$). Sphärisch: Winkelsumme $- (n - 2)\pi$ (bei $K = 1$). Euklidisch: null --- und deshalb keine Verbindung zwischen Winkeln und Größe. Gauß--Bonnet fasst alle drei zusammen (Aufgabe 6.33).
\end{mitnehmen}
\end{loesungkarte}

%%%%%%%%%% 6.23
\stupserzwei{gruen}{Winkelsumme $120^{\circ} = \frac{2\pi}3$; Defekt $\frac\pi3$. Mit $R = 3$: mal $R^{2} = 9$.}{22pt}
\begin{loesungkarte}{gruen}{6.23}{Ein hyperbolisches Dreieck}
\derweg{Defekt ausrechnen, mit $R^{2}$ skalieren, durch das Maximum $\pi$ teilen.}
\schritt{(a)}
Winkelsumme $120^{\circ} = \frac{2\pi}3$. Fläche $= \pi - \frac{2\pi}3 = \frac\pi3 = 1{,}047$.
\schritt{(b)}
Bei $K = -\frac1{R^{2}}$: Fläche $= R^{2}\cdot\text{Defekt} = 9\cdot\frac\pi3 = 3\pi = 9{,}42$. Dasselbe Dreieck (dieselben Winkel) ist in der flacheren Ebene neunmal größer --- Winkel bestimmen die Größe.
\schritt{(c)}
Maximum (bei $K = -1$): $\pi$. Anteil $\frac{\pi/3}\pi = \frac13$. Bei $R = 3$ ist das Maximum $9\pi$, der Anteil derselbe.
\probe{Ein euklidisches Dreieck mit Winkeln $30^{\circ}, 40^{\circ}, 50^{\circ}$ gibt es nicht ($120^{\circ}\neq180^{\circ}$); ein sphärisches auch nicht ($< 180^{\circ}$). Nur hyperbolisch. \checkmark}
\begin{mitnehmen}\kastentext
Das Defizit zu $180^{\circ}$ ist die Fläche in Einheiten von $R^{2}$. Für ein Dreieck von $1$ km auf einer hyperbolischen Ebene mit $R = 10^{6}$ km wäre das Defizit $10^{-12}$ --- unmessbar. Gauß vermaß das Dreieck Brocken--Hohenhagen--Inselsberg (bis $107$ km) und fand im Rahmen der Messgenauigkeit $180^{\circ}$: kein Beweis für Euklid, nur eine untere Schranke für $R$.
\end{mitnehmen}
\end{loesungkarte}

%%%%%%%%%% 6.24
\stupserzwei{akzent}{(b) $10^{\circ}$ auf $60^{\circ}$ N: Breitenkreis $556{,}0$ km, Großkreis $555{,}3$ km. Der Unterschied ist winzig, aber da.}{22pt}
\begin{loesungkarte}{akzent}{6.24}{Der Breitenkreis ist kein Umweg --- oder doch?}
\derweg{Breitenkreisbogen $= R\cos\varphi\cdot\Delta\lambda$; Großkreis über den sphärischen Kosinussatz; Geodäte heißt Ebene durch den Mittelpunkt.}
\schritt{(a)}
Breitenkreis bei $60^{\circ}$ N hat Radius $R\cos60^{\circ} = \frac R2$; der halbe Umfang ist $\pi\cdot\frac R2 = 10\,008$ km. Über den Pol: von $60^{\circ}$ N zum Pol sind $30^{\circ}$, hin und zurück $60^{\circ} = \frac\pi3$: $\frac\pi3R = 6672$ km. Der Breitenkreis ist $50\,\%$ länger.
\schritt{(b)}
Breitenkreis: $R\cos60^{\circ}\cdot\frac{10\pi}{180} = 556{,}0$ km. Großkreis: $\cos d = \sin^{2}60^{\circ} + \cos^{2}60^{\circ}\cos10^{\circ} = 0{,}75 + 0{,}25\cdot0{,}98481 = 0{,}99620$, $d = 4{,}994^{\circ}$, Strecke $555{,}3$ km. Unterschied $0{,}7$ km --- für kurze Strecken ist der Breitenkreis fast optimal, für lange nicht.
\schritt{(c)}
Geodäten auf der Kugel sind Großkreise: Schnitte mit Ebenen durch den Mittelpunkt. Die Äquatorebene geht durch den Mittelpunkt; die Ebene des Breitenkreises $60^{\circ}$ liegt $R\sin60^{\circ}$ darüber. Ein Breitenkreis hat Krümmungsradius $R\cos\varphi < R$ --- er biegt stärker als nötig, nach Norden hin.
\schritt{(d)}
Frankfurt und New York liegen auf $50^{\circ}$ und $41^{\circ}$ N; der Großkreis zwischen ihnen wölbt sich nach Norden, über Irland, Südgrönland, Neufundland --- bis $55^{\circ}$ N. Der Startkurs ist $295^{\circ}$ (Aufgabe 6.26), nicht $270^{\circ}$ (West), und ändert sich unterwegs.
\begin{fehler}\kastentext
Auf der Mercator-Karte sind Breitenkreise Geraden und Großkreise Bögen --- die Karte lügt über Längen (Aufgabe 6.30). Ein Breitenkreis ist nur am Äquator eine Geodäte. Die Intuition \enquote{geradeaus nach Westen ist der kürzeste Weg nach Westen} stimmt auf der Kugel nicht.
\end{fehler}
\begin{mitnehmen}\kastentext
Geodäte: lokal kürzeste Verbindung, Krümmung nur senkrecht zur Fläche. Auf der Kugel Großkreise, auf dem Zylinder Schraubenlinien (Aufgabe 6.29), in der hyperbolischen Ebene die Kreisbögen des Poincaré-Modells. Breitenkreise sind \enquote{Kreise um den Pol}, keine Geraden --- so wenig wie ein Kreis in der Ebene.
\end{mitnehmen}
\end{loesungkarte}

%%%%%%%%%% 6.25
\stupserzwei{gruen}{(c) Fläche $\approx\frac{\sqrt3}4\cdot100^{2} = 4330$ km$^{2}$, Exzess $= \frac A{R^{2}} = 1{,}07\cdot10^{-4}$ rad $= 22$ Bogensekunden.}{22pt}
\begin{loesungkarte}{gruen}{6.25}{Ein Dreieck auf der Erde}
\derweg{Winkel aus der Geometrie der Meridiane, Exzessformel, Vergleich; für das kleine Dreieck die euklidische Fläche als Näherung.}
\schritt{(a)}
Am Pol: Die Meridiane $0^{\circ}$ und $90^{\circ}$ bilden den Winkel $90^{\circ}$. Am Äquator trifft jeder Meridian senkrecht: $90^{\circ}$, $90^{\circ}$. Winkelsumme $270^{\circ}$, Exzess $90^{\circ} = \frac\pi2$. Fläche $= R^{2}\cdot\frac\pi2 = 6371^{2}\cdot1{,}5708 = 6{,}376\cdot10^{7}$ km$^{2}$.
\schritt{(b)}
Kugeloberfläche $4\pi R^{2} = 5{,}101\cdot10^{8}$ km$^{2}$, ein Achtel: $6{,}376\cdot10^{7}$. \checkmark\ Das Dreieck ist ein Oktant.
\schritt{(c)}
Für ein kleines Dreieck ist die Fläche fast euklidisch: gleichseitig mit $100$ km: $A = \frac{\sqrt3}4\cdot10^{4} = 4330$ km$^{2}$. Exzess $= \frac A{R^{2}} = \frac{4330}{4{,}059\cdot10^{7}} = 1{,}07\cdot10^{-4}$ rad $= 0{,}0061^{\circ} = 22''$. Mit Theodoliten (Genauigkeit $1''$) messbar; Gauß' Hannoversche Landesvermessung sah den Exzess.
\probe{Drei Oktanten zu $270^{\circ}$ plus\dots: Acht Oktanten mit je Exzess $\frac\pi2$ geben $4\pi = \frac{4\pi R^{2}}{R^{2}}$. \checkmark}
\begin{mitnehmen}\kastentext
Sphärischer Exzess $= \frac{\text{Fläche}}{R^{2}}$ --- für die Erde heißt das $22''$ pro $4330$ km$^{2}$, für ein Fußballfeld $10^{-10}$ Bogensekunden. Die Erde ist lokal flach, aber nicht global: Jede Landesvermessung muss den Exzess berücksichtigen, jede Stadtvermessung darf ihn ignorieren.
\end{mitnehmen}
\end{loesungkarte}

%%%%%%%%%% 6.26
\stupserzwei{gruen}{$\cos d = \sin50^{\circ}\sin40{,}7^{\circ} + \cos50^{\circ}\cos40{,}7^{\circ}\cos82{,}6^{\circ} = 0{,}4995 + 0{,}0628 = 0{,}5623$, $d = 55{,}78^{\circ}$. Kurs aus dem Kosinussatz für die Seite $NY$.}{22pt}
\begin{loesungkarte}{gruen}{6.26}{Frankfurt--New York}
\derweg{Sphärisches Dreieck Nordpol--Frankfurt--New York; Kosinussatz für die Distanz, noch einmal Kosinussatz für den Winkel in Frankfurt.}
\schritt{Dreieck}
Ecken $N$ (Pol), $F$, $Y$. Seiten vom Pol: $NF = 90^{\circ} - 50{,}0^{\circ} = 40{,}0^{\circ}$, $NY = 90^{\circ} - 40{,}7^{\circ} = 49{,}3^{\circ}$. Winkel am Pol: Längendifferenz $8{,}6^{\circ} + 74{,}0^{\circ} = 82{,}6^{\circ}$.
\schritt{Distanz}
Sphärischer Kosinussatz: $\cos d = \cos40^{\circ}\cos49{,}3^{\circ} + \sin40^{\circ}\sin49{,}3^{\circ}\cos82{,}6^{\circ} = 0{,}7660\cdot0{,}6521 + 0{,}6428\cdot0{,}7581\cdot0{,}1288 = 0{,}4995 + 0{,}0628 = 0{,}5623$. $d = 55{,}78^{\circ} = 0{,}9736$ rad. Strecke $R\cdot d = 6371\cdot0{,}9736 = 6203$ km.
\schritt{Startkurs}
Winkel $\alpha$ bei $F$ zwischen Meridian (nach Norden) und Großkreis: Kosinussatz für die Seite $NY$: $\cos49{,}3^{\circ} = \cos40^{\circ}\cos d + \sin40^{\circ}\sin d\cos\alpha$. $0{,}6521 = 0{,}7660\cdot0{,}5623 + 0{,}6428\cdot0{,}8270\cos\alpha$, $\cos\alpha = \frac{0{,}6521 - 0{,}4307}{0{,}5316} = 0{,}4165$, $\alpha = 65{,}4^{\circ}$. Der Kurs ist $65{,}4^{\circ}$ westlich von Nord: $360^{\circ} - 65{,}4^{\circ} = 294{,}6^{\circ}$ (Westnordwest).
\schritt{Vergleich}
Breitenkreisweg (ungefähr, auf $45^{\circ}$ N): $R\cos45^{\circ}\cdot82{,}6^{\circ}$ in rad $= 6371\cdot0{,}7071\cdot1{,}4416 = 6494$ km --- $290$ km länger. Der nördlichste Punkt des Großkreises liegt bei etwa $55^{\circ}$ N.
\probe{Sinussatz: $\frac{\sin\alpha}{\sin49{,}3^{\circ}} = \frac{\sin82{,}6^{\circ}}{\sin d}$: $\sin\alpha = \frac{0{,}7581\cdot0{,}9917}{0{,}8270} = 0{,}909$, $\alpha = 65{,}4^{\circ}$ (oder $114{,}6^{\circ}$ --- der Kosinussatz entscheidet). \checkmark\ Flugzeit bei $850$ km/h: $7{,}3$ h. \checkmark}
\begin{mitnehmen}\kastentext
Sphärischer Kosinussatz $\cos c = \cos a\cos b + \sin a\sin b\cos\gamma$: für kleine Dreiecke wird daraus (Taylor) der ebene Kosinussatz. Die Navigation rechnet seit Jahrhunderten damit; GPS-Empfänger tun es heute mit dem Ellipsoid statt der Kugel, der Unterschied liegt bei $0{,}3\,\%$.
\end{mitnehmen}
\end{loesungkarte}

%%%%%%%%%% 6.27
\stupserzwei{tief}{(b) $\kappa(0) = 2$, $\rho = \frac12$; $\kappa(1) = \frac2{5^{3/2}} = 0{,}179$, $\rho = 5{,}59$. (d) $\int_0^1\sqrt{1 + \sinh^{2}x}\,dx = \int_0^1\cosh x\,dx = \sinh1$.}{22pt}
\begin{loesungkarte}{tief}{6.27}{Krümmung eines Graphen}
\derweg{Graph als Kurve $(x, f(x))$ in die Parameterformel; dann einsetzen; die Bogenlänge der Kettenlinie als Geschenk.}
\schritt{(a)}
$\gamma(x) = (x, f(x))$, $\gamma' = (1, f')$, $\gamma'' = (0, f'')$. $\kappa = \frac{|x'y'' - y'x''|}{(x'^{2} + y'^{2})^{3/2}} = \frac{|1\cdot f'' - f'\cdot0|}{(1 + f'^{2})^{3/2}} = \frac{|f''|}{(1 + f'^{2})^{3/2}}$. $\blacksquare$
\schritt{(b)}
$f = x^{2}$, $f' = 2x$, $f'' = 2$. Bei $0$: $\kappa = 2$, $\rho = \frac12$ --- der Schmiegkreis hat Mittelpunkt $(0, \frac12)$ und berührt die Parabel im Scheitel von innen. Bei $1$: $\kappa = \frac2{(1 + 4)^{3/2}} = \frac2{11{,}18} = 0{,}179$, $\rho = 5{,}59$.
\schritt{(c)}
$f''$ ist die Änderung der \emph{Steigung} pro Einheit $x$; die Krümmung ist die Änderung des \emph{Winkels} pro Einheit \emph{Bogenlänge}. Bei steiler Tangente ändert sich der Winkel langsamer als die Steigung (Faktor $\frac1{1 + f'^{2}}$), und pro $x$-Einheit läuft man mehr Bogen (Faktor $\sqrt{1 + f'^{2}}$). Zusammen $(1 + f'^{2})^{3/2}$. Die Parabel hat überall $f'' = 2$, aber weit draußen ist sie fast gerade.
\schritt{(d)}
$y = \cosh x$, $y' = \sinh x$, $1 + y'^{2} = \cosh^{2}x$. $L = \int_0^1\cosh x\,dx = \sinh1 = 1{,}1752$. Die Kettenlinie ist die einzige Kurve, bei der $\sqrt{1 + y'^{2}} = y$ --- Bogenlänge gleich Fläche unter der Kurve.
\probe{Krümmung der Kettenlinie: $\kappa = \frac{\cosh x}{\cosh^{3}x} = \frac1{\cosh^{2}x}$, bei $0$ gleich $1$: Schmiegkreisradius $1$, Mittelpunkt $(0, 2)$. \checkmark}
\begin{mitnehmen}\kastentext
$\kappa = |f''|$ nur dort, wo $f' = 0$. Sonst ist die Krümmung kleiner --- deshalb kann eine Funktion mit $f'' = $ const (Parabel) im Unendlichen beliebig gerade werden, während der Kreis mit $\kappa = $ const sich schließt.
\end{mitnehmen}
\end{loesungkarte}

%%%%%%%%%% 6.28
\stupserzwei{gruen}{Bei $t = 0$: Nenner $4^{3/2} = 8$, $\kappa = \frac68 = \frac34$. Bei $t = \frac\pi2$: Nenner $9^{3/2} = 27$, $\kappa = \frac6{27} = \frac29$.}{22pt}
\begin{loesungkarte}{gruen}{6.28}{Die Krümmung der Ellipse}
\derweg{Krümmungsformel für Parameterkurven; Zähler konstant; Nenner in den Scheiteln.}
\schritt{Ableitungen}
$\gamma' = (-3\sin t, 2\cos t)$, $\gamma'' = (-3\cos t, -2\sin t)$. $x'y'' - y'x'' = 6\sin^{2}t + 6\cos^{2}t = 6$ --- konstant (die Flächengeschwindigkeit, zweites Keplersches Gesetz).
\schritt{Krümmung}
$\kappa(t) = \frac6{(9\sin^{2}t + 4\cos^{2}t)^{3/2}}$. Hauptscheitel $t = 0$, Punkt $(3, 0)$: $\kappa = \frac6{4^{3/2}} = \frac68 = \frac34$, $\rho = \frac43$. Nebenscheitel $t = \frac\pi2$, Punkt $(0, 2)$: $\kappa = \frac6{9^{3/2}} = \frac29$, $\rho = \frac92$.
\schritt{Allgemein}
$\rho = \frac{b^{2}}a$ im Hauptscheitel, $\frac{a^{2}}b$ im Nebenscheitel: $\frac43$ und $\frac92$. \checkmark\ Die Ellipse ist an den Enden der langen Achse am stärksten gekrümmt --- der Schmiegkreis dort ist klein; an den Enden der kurzen Achse ist der Schmiegkreis größer als die Ellipse selbst.
\probe{Kreis $a = b = r$: $\rho = r$ in beiden Scheiteln. \checkmark\ Krümmung im Hauptscheitel über die Graphenformel: nahe $(3, 0)$ ist $x = 3\sqrt{1 - y^{2}/4}\approx3 - \frac{3y^{2}}8$, also $|x''| = \frac34$ bei waagerechter Tangente. \checkmark}
\begin{mitnehmen}\kastentext
Der Schmiegkreis ist der Kreis, der die Kurve bis zur zweiten Ordnung berührt. Die Mittelpunkte aller Schmiegkreise bilden die Evolute --- bei der Ellipse eine Astroide. Bei Eisenbahnkurven schaltet man zwischen Gerade und Kreisbogen eine Klothoide, damit die Krümmung stetig wächst; ein Sprung in $\kappa$ ist ein Ruck.
\end{mitnehmen}
\end{loesungkarte}

%%%%%%%%%% 6.29
\stupserzwei{gruen}{$\gamma'\times\gamma'' = (ab\sin t, -ab\cos t, a^{2})$, Betrag $a\sqrt{a^{2} + b^{2}}$; $\gamma''' = (a\sin t, -a\cos t, 0)$; $(\gamma'\times\gamma'')\cdot\gamma''' = a^{2}b$.}{22pt}
\begin{loesungkarte}{gruen}{6.29}{Die Schraubenlinie}
\derweg{Drei Ableitungen, Kreuzprodukt, die beiden Formeln; dann $b\to0$ und $a\to0$.}
\schritt{Ableitungen}
$\gamma' = (-a\sin t, a\cos t, b)$, $|\gamma'| = \sqrt{a^{2} + b^{2}}$ (konstant: die Kurve wird mit konstanter Geschwindigkeit durchlaufen). $\gamma'' = (-a\cos t, -a\sin t, 0)$, $\gamma''' = (a\sin t, -a\cos t, 0)$.
\schritt{Krümmung}
$\gamma'\times\gamma'' = (a\cos t\cdot0 - b(-a\sin t),\ b(-a\cos t) - (-a\sin t)\cdot0,\ a^{2}\sin^{2}t + a^{2}\cos^{2}t) = (ab\sin t, -ab\cos t, a^{2})$. Betrag $\sqrt{a^{2}b^{2} + a^{4}} = a\sqrt{a^{2} + b^{2}}$. $\kappa = \frac{a\sqrt{a^{2} + b^{2}}}{(a^{2} + b^{2})^{3/2}} = \frac a{a^{2} + b^{2}}$.
\schritt{Torsion}
$(\gamma'\times\gamma'')\cdot\gamma''' = a^{2}b\sin^{2}t + a^{2}b\cos^{2}t + 0 = a^{2}b$. $\tau = \frac{a^{2}b}{a^{2}(a^{2} + b^{2})} = \frac b{a^{2} + b^{2}}$.
\schritt{Grenzfälle}
$b = 0$: Kreis mit Radius $a$, $\kappa = \frac1a$, $\tau = 0$ --- eben. $a\to0$: Gerade, $\kappa\to0$, $\tau\to0$ (bei $a = 0$ ist die Torsion nicht definiert). $b < 0$: Linksschraube, $\tau < 0$. Für $a = b$: $\kappa = \tau = \frac1{2a}$.
\probe{$\kappa^{2} + \tau^{2} = \frac{a^{2} + b^{2}}{(a^{2} + b^{2})^{2}} = \frac1{a^{2} + b^{2}}$ --- die Schraubenlinie ist eine \enquote{Kreislinie} vom Radius $\sqrt{a^{2} + b^{2}}$ in einem gewissen Sinn. \checkmark}
\begin{mitnehmen}\kastentext
Hauptsatz der Kurventheorie: $\kappa(s)$ und $\tau(s)$ bestimmen die Kurve bis auf Bewegung. Konstante $\kappa, \tau$ geben genau die Schraubenlinien (einschließlich Kreis und Gerade). DNA, Sprungfedern, Schiffsschrauben: Die Natur mag konstante Krümmung und Torsion.
\end{mitnehmen}
\end{loesungkarte}

%%%%%%%%%% 6.30
\stupserzwei{akzent}{(c) Winkeltreu und flächentreu zugleich wäre längentreu --- und das verbietet $K\neq0$. Grönland ($2{,}2$ Mio.\ km$^{2}$) erscheint auf Mercator so groß wie Afrika ($30$ Mio.).}{22pt}
\begin{loesungkarte}{akzent}{6.30}{Warum keine Landkarte stimmt}
\derweg{Flächenelement aus der Fundamentalform; Theorema egregium: Isometrien erhalten $K$; daraus alle Unmöglichkeiten.}
\schritt{(a)}
Das Flächenelement ist $\sqrt{EG - F^{2}}\,d\theta\,d\varphi = R^{2}\sin\theta\,d\theta\,d\varphi$ (Aufgabe 6.31). Der Faktor $\sin\theta$ fehlt: Nahe den Polen sind die Koordinatenmaschen schmal. Richtig: $\int_0^{2\pi}\int_0^{\pi}R^{2}\sin\theta\,d\theta\,d\varphi = 2\pi R^{2}\cdot2 = 4\pi R^{2}$.
\schritt{(b)}
Theorema egregium: Die Gauß-Krümmung hängt nur von der ersten Fundamentalform ab, also bleibt sie bei längentreuen Abbildungen (Isometrien) erhalten. Zylinder: $K = 0$ (eine Hauptkrümmung null), Ebene: $K = 0$ --- Abwickeln möglich, und man sieht es: Ein Blatt Papier lässt sich rollen. Kugel: $K = \frac1{R^{2}} > 0$, Ebene $K = 0$ --- keine Isometrie, nicht einmal lokal. Eine Orangenschale lässt sich nicht flach drücken, ohne zu reißen.
\schritt{(c)}
Winkeltreu und flächentreu heißt: Die lineare Näherung ist eine Ähnlichkeit (winkeltreu) mit Determinante $1$ (flächentreu), also eine Isometrie --- unmöglich nach (b). Mercator ist winkeltreu (Kurse sind Geraden, deshalb für die Seefahrt gemacht) und streckt dafür Flächen um $\frac1{\cos^{2}\varphi}$: bei $70^{\circ}$ N um den Faktor $8{,}5$. Grönland sieht aus wie Afrika, ist aber $14$-mal kleiner.
\schritt{(d)}
Nein. Längentreu auf einem noch so kleinen Stück hieße $K = 0$ dort, und $K = \frac1{R^{2}}$ überall. Jede Karte verzerrt --- man kann nur wählen, was: Winkel (Mercator), Flächen (Peters, Lambert), Abstände von einem Punkt (azimutal). Das ist auch der Grund für die Sektoren eines Globus-Bastelbogens: viele schmale Streifen, jeder fast flach.
\begin{fehler}\kastentext
Zwei Fehler, eine Ursache: Das Flächenelement $\sqrt{EG - F^{2}}$ ist die Verzerrung der Koordinatenkarte, und $K\neq0$ heißt, dass keine Karte sie beseitigt. Wer $d\theta\,d\varphi$ integriert, behandelt die Karte, als wäre sie die Fläche.
\end{fehler}
\begin{mitnehmen}\kastentext
Gauß' \enquote{Theorema egregium} (1827): Krümmung ist eine innere Größe --- ein zweidimensionales Wesen auf der Fläche kann sie messen, ohne nach außen zu schauen (Winkelsumme kleiner Dreiecke, Umfang kleiner Kreise). Das ist der Gedanke, den Riemann zu $n$ Dimensionen und Einstein zur Gravitation gemacht hat.
\end{mitnehmen}
\end{loesungkarte}

%%%%%%%%%% 6.31
\stupserzwei{gruen}{$X_\theta$ zeigt nach Süden, $X_\varphi$ nach Osten; beide stehen senkrecht ($F = 0$). $E = R^{2}$, $G = R^{2}\sin^{2}\theta$ --- ein Grad Länge ist $\sin\theta$-mal so lang wie ein Grad Breite.}{22pt}
\begin{loesungkarte}{gruen}{6.31}{Die Kugel und ihre Oberfläche}
\derweg{Parametrisierung ableiten, drei Skalarprodukte, Flächenelement integrieren.}
\schritt{Parametrisierung}
$X(\theta, \varphi) = R(\sin\theta\cos\varphi, \sin\theta\sin\varphi, \cos\theta)$, $\theta\in(0, \pi)$ Poldistanz, $\varphi\in[0, 2\pi)$ Länge.
\schritt{Tangentenvektoren}
$X_\theta = R(\cos\theta\cos\varphi, \cos\theta\sin\varphi, -\sin\theta)$ (nach Süden), $X_\varphi = R(-\sin\theta\sin\varphi, \sin\theta\cos\varphi, 0)$ (nach Osten).
\schritt{Fundamentalform}
$E = X_\theta\cdot X_\theta = R^{2}(\cos^{2}\theta + \sin^{2}\theta) = R^{2}$. $F = X_\theta\cdot X_\varphi = R^{2}(-\cos\theta\sin\theta\cos\varphi\sin\varphi + \cos\theta\sin\theta\sin\varphi\cos\varphi) = 0$. $G = X_\varphi\cdot X_\varphi = R^{2}\sin^{2}\theta$. \\
$ds^{2} = R^{2}\,d\theta^{2} + R^{2}\sin^{2}\theta\,d\varphi^{2}$: Meridiane und Breitenkreise stehen senkrecht ($F = 0$); ein Grad Länge ist $\sin\theta$-mal so lang wie ein Grad Breite.
\schritt{Oberfläche}
$dA = \sqrt{EG - F^{2}}\,d\theta\,d\varphi = R^{2}\sin\theta\,d\theta\,d\varphi$. $A = \int_0^{2\pi}\!\!\int_0^{\pi}R^{2}\sin\theta\,d\theta\,d\varphi = 2\pi R^{2}[-\cos\theta]_0^{\pi} = 4\pi R^{2}$.
\probe{Kugelzone zwischen $\theta_1$ und $\theta_2$: $2\pi R^{2}(\cos\theta_1 - \cos\theta_2)$, hängt nur von der Höhe $h = R(\cos\theta_1 - \cos\theta_2)$ ab: $2\pi Rh$ --- Archimedes' Satz. \checkmark\ Breitenkreislänge bei $\theta$: $\int_0^{2\pi}\sqrt G\,d\varphi = 2\pi R\sin\theta$. \checkmark}
\begin{mitnehmen}\kastentext
Die erste Fundamentalform $ds^{2} = E\,du^{2} + 2F\,du\,dv + G\,dv^{2}$ ist alles, was ein Flächenbewohner kennt: Längen ($\int\sqrt{E\dot u^{2} + \dots}$), Winkel ($F = 0$ heißt orthogonale Koordinaten), Flächen ($\sqrt{EG - F^{2}}$) --- und nach Gauß auch die Krümmung.
\end{mitnehmen}
\end{loesungkarte}

%%%%%%%%%% 6.32
\stupserzwei{gruen}{Normale $N = (\cos u\cos v, \cos u\sin v, \sin u)$; $L = r$, $M = 0$, $N_{\text{II}} = (R + r\cos u)\cos u$; $E = r^{2}$, $F = 0$, $G = (R + r\cos u)^{2}$.}{22pt}
\begin{loesungkarte}{gruen}{6.32}{Die Krümmung des Torus}
\derweg{Erste und zweite Fundamentalform, $K = \frac{LN - M^{2}}{EG - F^{2}}$, Vorzeichen von $\cos u$; dann das Integral.}
\schritt{Erste Fundamentalform}
$X_u = (-r\sin u\cos v, -r\sin u\sin v, r\cos u)$, $X_v = (-(R + r\cos u)\sin v, (R + r\cos u)\cos v, 0)$. $E = r^{2}$, $F = 0$, $G = (R + r\cos u)^{2}$.
\schritt{Normale und zweite Fundamentalform}
$N = \frac{X_u\times X_v}{|X_u\times X_v|} = (\cos u\cos v, \cos u\sin v, \sin u)$ (nach außen). $X_{uu} = (-r\cos u\cos v, -r\cos u\sin v, -r\sin u)$, $X_{uv} = (r\sin u\sin v, -r\sin u\cos v, 0)$, $X_{vv} = (-(R + r\cos u)\cos v, -(R + r\cos u)\sin v, 0)$. \\
$L = X_{uu}\cdot N = -r$, $M = X_{uv}\cdot N = 0$, $\mathcal N = X_{vv}\cdot N = -(R + r\cos u)\cos u$. (Vorzeichen nach Konvention; das Produkt $L\mathcal N$ ist davon unabhängig.)
\schritt{Gauß-Krümmung}
$K = \frac{L\mathcal N - M^{2}}{EG - F^{2}} = \frac{r(R + r\cos u)\cos u}{r^{2}(R + r\cos u)^{2}} = \frac{\cos u}{r(R + r\cos u)}$.
\schritt{Typen}
$\cos u > 0$ (Außenseite, $|u| < \frac\pi2$): $K > 0$, elliptisch --- wie eine Kugel. $\cos u < 0$ (Innenseite): $K < 0$, hyperbolisch --- Sattel. $\cos u = 0$ (oberster und unterster Kreis): $K = 0$, parabolisch. Maximum $K = \frac1{r(R + r)}$ ganz außen, Minimum $\frac{-1}{r(R - r)}$ ganz innen --- innen ist die Krümmung betragsmäßig größer.
\schritt{Integral}
$dA = \sqrt{EG}\,du\,dv = r(R + r\cos u)\,du\,dv$. $\int_TK\,dA = \int_0^{2\pi}\!\!\int_0^{2\pi}\cos u\,du\,dv = 2\pi\cdot0 = 0$. Gauß--Bonnet: $2\pi\chi(T) = 0$. \checkmark
\probe{Grenzfall $R\to\infty$ bei festem $r$: $K\to\frac{\cos u}{rR}\to0$ --- der Torus wird zum Zylinder. \checkmark\ Außen bei $u = 0$: Hauptkrümmungen $\frac1r$ (Meridian) und $\frac1{R + r}$ (Breitenkreis), Produkt stimmt. \checkmark}
\begin{mitnehmen}\kastentext
$K = \kappa_1\kappa_2$: Produkt der Hauptkrümmungen. Positiv, wenn die Fläche in beide Richtungen gleich gekrümmt ist (Kuppe oder Mulde), negativ beim Sattel, null, wenn eine Richtung gerade ist. Der Torus zeigt alle drei --- und sein Gesamtintegral ist null, egal wie dick er ist.
\end{mitnehmen}
\end{loesungkarte}

%%%%%%%%%% 6.33
\stupserzwei{tief}{(c) $E = \frac{12\cdot5 + 20\cdot6}2 = 90$, $V = \frac{2E}3 = 60$ (drei Kanten pro Ecke), $F = 32$: $60 - 90 + 32 = 2$.}{22pt}
\begin{loesungkarte}{tief}{6.33}{Gauß--Bonnet konkret}
\derweg{Beide Seiten der Formel ausrechnen, wo man kann; die Topologie zählt Ecken, Kanten, Flächen.}
\schritt{(a)}
$K = \frac1{R^{2}}$, $A = 4\pi R^{2}$: $\int K\,dA = 4\pi$. $\chi(\text{Kugel}) = 2$, $2\pi\chi = 4\pi$. \checkmark\ Unabhängig von $R$: Eine kleine Kugel ist stark gekrümmt, aber klein.
\schritt{(b)}
Dreieck mit drei rechten Winkeln (Oktant): Fläche $\frac{\pi R^{2}}2$, $\int K\,dA = \frac\pi2$. Gauß--Bonnet für geodätische Dreiecke: $\int K\,dA = \alpha + \beta + \gamma - \pi = \frac{3\pi}2 - \pi = \frac\pi2$. \checkmark\ Das ist der sphärische Exzess (Aufgabe 6.25).
\schritt{(c)}
$F = 32$. Kanten: jede Fläche hat $5$ bzw.\ $6$, jede Kante gehört zu zwei Flächen: $E = \frac{60 + 120}2 = 90$. Ecken: an jeder treffen sich $3$ Kanten: $V = \frac{2\cdot90}3 = 60$. $V - E + F = 60 - 90 + 32 = 2$. \checkmark\ Und mit $12$ Fünfecken geht es immer: Für ein Polyeder aus Fünf- und Sechsecken mit drei Kanten pro Ecke folgt aus Euler $F_5 = 12$, egal wie viele Sechsecke.
\schritt{(d)}
Eine geschlossene Fläche ohne Loch ist topologisch eine Kugel, $\chi = 2$, also $\int K\,dA = 4\pi > 0$. Irgendwo muss $K > 0$ sein --- egal wie viele Dellen (negativ gekrümmte Stellen) man hineindrückt, die Kuppen gleichen sie aus. Für den Torus ($\chi = 0$) ist das Integral null (Aufgabe 6.32), für eine Brezel mit $g$ Löchern $4\pi(1 - g) < 0$: Flächen mit vielen Löchern sind überwiegend sattelförmig.
\probe{Würfel als \enquote{Kugel mit Ecken}: $K$ ist auf den Flächen null, konzentriert in den $8$ Ecken mit je Defekt $2\pi - 3\cdot\frac\pi2 = \frac\pi2$: $8\cdot\frac\pi2 = 4\pi$. \checkmark\ Descartes' Satz vom Winkeldefekt ist Gauß--Bonnet für Polyeder.}
\begin{mitnehmen}\kastentext
$\int_MK\,dA = 2\pi\chi(M)$: Links Analysis (Krümmung, punktweise), rechts Topologie (eine ganze Zahl). Verbiegen ändert $K$ überall, aber nicht das Integral. Das ist die erste der großen Index-Formeln --- die Vorfahrin von Riemann--Roch, Atiyah--Singer und dem Chern-Gauß-Bonnet-Satz.
\end{mitnehmen}
\end{loesungkarte}

%%%%%%%%%% 6.34
\stupserzwei{gruen}{Torus: $\chi = 0$. Das Feld \enquote{immer nach Osten} ($X_v$) hat keine Nullstelle --- es kämmt den Torus.}{22pt}
\begin{loesungkarte}{gruen}{6.34}{Der Igelsatz}
\derweg{Index einer Nullstelle definieren; Poincaré--Hopf verbindet die Indexsumme mit $\chi$; Gauß--Bonnet liefert $\chi$ aus der Krümmung.}
\schritt{Index}
Hat ein Tangentialvektorfeld bei $P$ eine Nullstelle, so umläuft man $P$ auf einem kleinen Kreis und zählt, wie oft sich der Feldvektor dabei dreht: Quelle, Senke, Wirbel haben Index $+1$, ein Sattel $-1$.
\schritt{Poincaré--Hopf}
Für jedes Vektorfeld mit endlich vielen Nullstellen auf einer geschlossenen Fläche $M$: $\sum\text{Index} = \chi(M)$. Auf der Kugel $\chi = 2$: Die Indexsumme ist $2\neq0$, also gibt es mindestens eine Nullstelle. Ein nirgends verschwindendes Feld hätte Indexsumme $0$. $\blacksquare$
\schritt{Zusammenhang mit Gauß--Bonnet}
Gauß--Bonnet sagt $\int K\,dA = 2\pi\chi$; Poincaré--Hopf sagt $\sum\text{Index} = \chi$. Beide berechnen dieselbe topologische Zahl, einmal aus der Krümmung, einmal aus einem Vektorfeld. Der Beweis von Poincaré--Hopf benutzt Gauß--Bonnet: Man integriert $K$ über die Fläche ohne kleine Scheiben um die Nullstellen; die Randterme sind die Indizes.
\schritt{Beispiele}
Das Feld \enquote{nach Norden} auf der Erde: Nullstellen an beiden Polen (Senke am Nordpol, Quelle am Südpol), Indizes $+1 + 1 = 2$. \checkmark\ Das Feld \enquote{nach Osten}: Wirbel an beiden Polen, $+1 + 1 = 2$. \checkmark\ Der Torus: $\chi = 0$; das Feld $X_v$ (entlang der Breitenkreise) hat keine Nullstelle. Der Torus lässt sich kämmen --- ein Donut mit Streuseln in Umlaufrichtung.
\probe{Ein Feld auf der Kugel mit nur \emph{einer} Nullstelle? Ja: stereographisch das konstante Feld der Ebene übertragen; am Nordpol entsteht eine Nullstelle mit Index $2$ (ein \enquote{Dipol}). \checkmark}
\begin{mitnehmen}\kastentext
\enquote{Auf der Erde ist immer irgendwo Windstille} --- ein Satz der Topologie, nicht der Meteorologie. Dieselbe Mathematik: Ein Wirbelsturm auf der Kugel braucht ein Gegenstück; ein Magnetfeld auf einer Kugel hat Pole; stetige Abbildungen der Kugel auf sich haben Fixpunkte oder Antipodenpunkte.
\end{mitnehmen}
\end{loesungkarte}

%%%%%%%%%% 6.35
\stupserzwei{tief}{$\frac{2E}q - E + \frac{2E}p = 2$, also $E\bigl(\frac2p + \frac2q - 1\bigr) = 2 > 0$, also $\frac1p + \frac1q > \frac12$.}{22pt}
\begin{loesungkarte}{tief}{6.35}{Die fünf platonischen Körper}
\derweg{Doppeltes Abzählen der Kanten, Euler einsetzen, die Ungleichung in ganzen Zahlen lösen.}
\schritt{(a)}
Jede der $F$ Flächen hat $p$ Kanten, jede Kante liegt in $2$ Flächen: $pF = 2E$, $F = \frac{2E}p$. An jeder der $V$ Ecken enden $q$ Kanten, jede Kante hat $2$ Enden: $qV = 2E$, $V = \frac{2E}q$. Euler: $\frac{2E}q - E + \frac{2E}p = 2$.
\schritt{(b)}
$E\bigl(\frac2p + \frac2q - 1\bigr) = 2$. Da $E > 0$: $\frac2p + \frac2q > 1$, also $\frac1p + \frac1q > \frac12$. Mit $p, q\ge3$: $p = 3$: $\frac1q > \frac16$, $q\in\{3, 4, 5\}$. $p = 4$: $\frac1q > \frac14$, $q = 3$. $p = 5$: $\frac1q > \frac3{10}$, $q = 3$. $p\ge6$: $\frac1q > \frac13$, unmöglich. Fünf Lösungen.
\schritt{(c)}
$E = \frac2{2/p + 2/q - 1}$, dann $V, F$:
\begin{center}\small
\begin{tabular}{cc|ccc|l}
$p$ & $q$ & $V$ & $E$ & $F$ & Körper\\ \hline
3 & 3 & 4 & 6 & 4 & Tetraeder\\
4 & 3 & 8 & 12 & 6 & Würfel\\
3 & 4 & 6 & 12 & 8 & Oktaeder\\
5 & 3 & 20 & 30 & 12 & Dodekaeder\\
3 & 5 & 12 & 30 & 20 & Ikosaeder
\end{tabular}
\end{center}
Vertauschen von $p$ und $q$ vertauscht $V$ und $F$: Dualität (Würfel--Oktaeder, Dodekaeder--Ikosaeder, Tetraeder mit sich selbst).
\schritt{(d)}
$\frac1p + \frac1q = \frac12$: $(3, 6)$, $(4, 4)$, $(6, 3)$. Dann wäre $E = \infty$ --- das sind die drei regelmäßigen Parkettierungen der Ebene: Dreiecke (sechs an jeder Ecke), Quadrate, Sechsecke. Und $\frac1p + \frac1q < \frac12$ liefert die Parkettierungen der hyperbolischen Ebene (Aufgabe 6.22: rechtwinklige Fünfecke, $(5, 4)$).
\probe{Dodekaeder: $12$ Fünfecke, $E = 30$, $V = 20$; $20 - 30 + 12 = 2$. \checkmark\ Winkeldefekt an jeder Ecke $360^{\circ} - 3\cdot108^{\circ} = 36^{\circ}$, mal $20$ Ecken: $720^{\circ} = 4\pi$. \checkmark}
\begin{mitnehmen}\kastentext
Euklids \emph{Elemente} enden mit diesem Satz (Buch XIII). Der Beweis über Euler ist moderner als Euklids über Winkelsummen und zeigt mehr: Die Ungleichung sortiert zugleich die sphärischen ($> \frac12$), euklidischen ($= \frac12$) und hyperbolischen ($< \frac12$) regelmäßigen Parkettierungen --- drei Geometrien in einer Ungleichung.
\end{mitnehmen}
\end{loesungkarte}

%%%%%%%%%% 6.36
\stupserzwei{gruen}{Läge auf $g_0$ ein dritter Punkt, so liegen zwei der drei auf derselben Seite des Lotfußpunkts von $P_0$ --- und der nähere von beiden hat zur Geraden durch $P_0$ und den ferneren einen kleineren Abstand als $P_0$ zu $g_0$.}{22pt}
\begin{loesungkarte}{gruen}{6.36}{Sylvester--Gallai}
\derweg{Extremalprinzip: unter endlich vielen Paaren (Punkt, Gerade) das mit kleinstem positivem Abstand; dann ein noch kleineres konstruieren, falls drei Punkte auf der Geraden liegen.}
\schritt{Kandidaten}
Betrachte alle Paare $(P, g)$: $g$ eine Gerade durch mindestens zwei der Punkte, $P$ ein Punkt der Menge nicht auf $g$. Da nicht alle kollinear sind, gibt es solche Paare; da die Menge endlich ist, endlich viele.
\schritt{Minimum}
Wähle $(P_0, g_0)$ mit kleinstem Abstand $d > 0$. Sei $F$ der Lotfußpunkt von $P_0$ auf $g_0$.
\schritt{Widerspruch}
Angenommen, auf $g_0$ liegen drei Punkte $A, B, C$ der Menge. Zwei davon liegen auf derselben Seite von $F$ (oder einer ist $F$) --- sagen wir $A$ und $B$ mit $FA\le FB$ (eventuell $A = F$). Betrachte die Gerade $h = P_0B$. Der Abstand von $A$ zu $h$ ist kleiner als $d$: Im Dreieck $P_0FB$ ist $A$ auf der Strecke $FB$, und die Höhe von $A$ auf $P_0B$ ist kürzer als die Höhe von $F$ auf $P_0B$, die wiederum kürzer als $P_0F = d$ ist (Höhe im rechtwinkligen Dreieck). Also hat $(A, h)$ einen kleineren positiven Abstand als $(P_0, g_0)$ --- Widerspruch zur Minimalität. Also liegen auf $g_0$ genau zwei Punkte. $\blacksquare$
\begin{center}
\begin{tikzpicture}[scale=0.9]
\draw[linie] (-2.5,0)--(3.5,0) node[right,grau,font=\small] {$g_0$};
\fill[tief] (0,1.6) circle (1.6pt) node[above,font=\small] {$P_0$};
\fill[grau] (0,0) circle (1.2pt) node[below,font=\small] {$F$};
\fill[tief] (1,0) circle (1.6pt) node[below,font=\small] {$A$};
\fill[tief] (2.8,0) circle (1.6pt) node[below,font=\small] {$B$};
\draw[akzent] (0,1.6)--(2.8,0) node[midway,above right,font=\scriptsize] {$h$};
\draw[dashed,grau] (0,1.6)--(0,0);
\draw[dashed,akzent] (1,0)--(1.44,0.78);
\end{tikzpicture}
\end{center}
\schritt{Unendlich viele Punkte}
Das Extremalprinzip braucht ein Minimum, und das gibt es bei unendlich vielen Paaren nicht unbedingt (Infimum $0$). Und der Satz ist dort falsch: Alle Gitterpunkte $\mathbb Z^{2}$ --- jede Verbindungsgerade enthält unendlich viele. Ebenso alle Punkte der Ebene.
\probe{Kleinstes Beispiel: $3$ Punkte, nicht kollinear --- jede Verbindungsgerade enthält genau zwei. \checkmark\ Reguläres $n$-Eck plus Mittelpunkt: Die Geraden durch den Mittelpunkt enthalten $3$ Punkte (bei geradem $n$), die Seiten je zwei. \checkmark}
\begin{mitnehmen}\kastentext
Extremalprinzip: Wähle das kleinste, größte, nächste Objekt --- und zeige, dass ein noch extremeres existiert, falls die Behauptung falsch wäre. Es braucht Endlichkeit (oder Kompaktheit), und es beweist Existenzsätze ohne Konstruktion. Kellys Beweis ist das Lehrbuchbeispiel; Erdős nahm ihn ins \enquote{Buch}.
\end{mitnehmen}
\end{loesungkarte}

%%%%%%%%%% 6.37
\stupserzwei{gruen}{Raum: kubisch-flächenzentrierte Zelle mit Kantenlänge $a = 2\sqrt2\,r$ (Flächendiagonale $= 4r$) enthält $4$ Kugeln: $\frac{4\cdot\frac43\pi r^{3}}{16\sqrt2\,r^{3}}$.}{22pt}
\begin{loesungkarte}{gruen}{6.37}{Dichteste Packungen}
\derweg{Elementarzelle wählen, Kugelanteil darin zählen, durch das Zellvolumen teilen.}
\schritt{Ebene: hexagonal}
Drei sich berührende Kreise vom Radius $r$; ihre Mittelpunkte bilden ein gleichseitiges Dreieck mit Seite $2r$, Fläche $\frac{\sqrt3}4(2r)^{2} = \sqrt3\,r^{2}$. Im Dreieck liegen drei Kreissektoren zu je $60^{\circ}$, zusammen ein halber Kreis: $\frac{\pi r^{2}}2$. Dichte $\frac{\pi r^{2}/2}{\sqrt3\,r^{2}} = \frac\pi{2\sqrt3} = 0{,}9069$.
\schritt{Raum: kubisch-flächenzentriert}
Kugeln an den Ecken und Flächenmitten eines Würfels mit Kante $a$; Nachbarn berühren sich entlang der Flächendiagonale: $a\sqrt2 = 4r$, $a = 2\sqrt2\,r$. Kugeln in der Zelle: $8$ Ecken zu je $\frac18$ plus $6$ Flächen zu je $\frac12$: $1 + 3 = 4$. Dichte $\frac{4\cdot\frac43\pi r^{3}}{(2\sqrt2\,r)^{3}} = \frac{\frac{16}3\pi}{16\sqrt2} = \frac\pi{3\sqrt2} = 0{,}7405$.
\schritt{Einordnung}
Ebene: Optimalität bewiesen von Thue (1910), streng von Fejes Tóth (1940). Raum: Keplers Vermutung (1611), bewiesen von Hales (1998, $250$ Seiten plus Computer; formal verifiziert 2014). Die hexagonale Schichtung (hcp) hat dieselbe Dichte --- es gibt unendlich viele optimale Packungen, alle aus den gleichen Schichten in verschiedener Stapelfolge.
\probe{Kubisch-primitiv (Kugeln nur an den Ecken, $a = 2r$): $\frac{\frac43\pi r^{3}}{8r^{3}} = \frac\pi6 = 0{,}524$. Kubisch-raumzentriert: $0{,}680$. Beide schlechter. \checkmark\ Zufällig geschüttete Kugeln: etwa $0{,}64$.}
\begin{mitnehmen}\kastentext
Elementarzelle: die kleinste Einheit, die durch Verschieben den ganzen Raum füllt; Dichte $=$ Volumenanteil darin. In Dimension $8$ und $24$ sind die optimalen Packungen bekannt ($E_8$, Leech-Gitter; Viazovska 2016), in Dimension $4$ bis $7$ nicht --- und die dichtesten Packungen in hohen Dimensionen sind die Grundlage fehlerkorrigierender Codes (Band 20).
\end{mitnehmen}
\end{loesungkarte}
